% Appendix A --- deferred proofs of section 3. Moved out of section 3 (Q-0291, the author's decision
% of 2026-10-04 to shorten sections 3 and 4 by moving material to appendices, as Q-0290 did for
% section 5; records/affine/move-sections-3-4-to-appendices*.py, three passes). The proofs are
% unchanged but for their headings, in main-text order: lem:radial-marginal,
% lem:isotropic-marginal-family, the sufficiency and uniqueness halves of
% thm:isotropic-classification, prop:isotropic-polar-reading, prop:isotropic-bridge (3.2);
% lem:positive-transform-origin, lem:isotropic-exponent-log, thm:origin-boundedness-d (3.3);
% lem:threshold-monotone, lem:threshold-limit, thm:filtration-strict, thm:no-cin-ray-in-the-plane
% (3.4). Every statement stays in section 3, which keeps the necessity half of
% thm:isotropic-classification: the reduction to the line.
%
% Q-0310 (the whole-paper review's notation flags; records/affine/
% REVIEW-paper-whole-fixes-paper-iso.md), in the proofs only: the contraction ratio of
% self-decomposability is lambda (it was c, which is the profile at the origin); the
% sphere-averaged profile is \bar k (K is the constant of lem:threshold-limit and the Bessel
% function); the unit vector of lem:isotropic-exponent-log's proof is \widehat\xi (theta is the
% Matern range); the shadow measure of thm:no-cin-ray-in-the-plane's proof is \bar\nu (m is
% (d+1)/2); "profile" is qualified. The gauge ratio c of the uniqueness proof matches the shared
% statement and is owed at the blueprint (records/affine/GATE-paper-iso.md).
%
% Q-0311 (the blind review records/affine/REVIEW-paper-iso.md, flag 3): the three proofs that were
% sketches (lem:isotropic-exponent-log, thm:origin-boundedness-d, lem:threshold-limit) are now the
% blueprint's proofs of record in full (blueprint/src/parts/03-isotropic.tex), transcribed with the
% paper's conventions: "Lemma~\ref" for a bare \ref, the cited facts by their paper names C1-C4,
% the unit vector \widehat\xi for the blueprint's theta (the Matern range), the radius r of the
% punctured balls in clause (3) for the blueprint's eta (the sphere variable), y for the dummy
% variable of clause (4), and C_k defined in clause (3), where the blueprint uses it for the delay
% profile's tail without a definition. No proof of the paper is a sketch now. The remark on
% 3/2 - sqrt 2 moved here from section 3.4, after its proof (flag 31).
%
% Q-0316 (pass 2, the blind register review REVIEW-register-paper-iso.md, Appendix A flags 1-25):
% the proof-zone flags were fixed at the blueprint's proofs of record, and every proof here is
% re-transcribed from them verbatim by records/affine/transcribe-proofs-q0316.py (the cited facts
% through \ledger, which macros.tex prints as C1-C5). Two proofs are new: those of
% prop:isotropic-radial-monotone and lem:delay-polar-profile (the witness's membership, flag 1).
% The remark after lem:threshold-monotone is paper prose (flag 22), shortened.
%
% Q-0328: the proof of record of prop:origin-lower-bound-d, verbatim (records/affine/
% transcribe-q0328-origin-lower-bound.py), after that of thm:origin-boundedness-d. Its register was
% brought to the Q-0316 pass's at the blueprint first (seven participial clauses and two semicolons
% made sentences, no step changed).
%
% Q-0406 (pass 6, the author's review of section 3; records/affine/REVIEW-author-paper-iso.md):
% A.1 is the proof of the one lemma on operators and kernels, in both directions, from the two
% proofs of record (the comment above it says how). Item 5 of the brief: a place of the line paper
% or the cone paper that section 2 restates is cited at the restatement, through \restated and
% \restatedp (macros.tex; the blueprint changed by the same markup, so the proofs stay verbatim);
% the proof of prop:isotropic-bridge drops a restatement that section 2 gives, and is pinned. Not
% touched: the proof of thm:filtration-strict, which Q-0381 is rewriting on another branch (its
% [B, Prop. 4.9(1)] and (4) are Proposition prop:line-filtration(1) and (3) and are owed).

% Q-0448 (the split into a paper and a supplement; the author's decision of 2026-10-13, option
% (b) of records/affine/reviews/VERIFY-round2-length-and-exposition.md, section 3): this file is
% Appendix A of the paper. It keeps the notation, the sources, the proofs of the classification,
% the statements that the main text names and that are proved in the supplement, and the
% statements that Lean did not check, with their proofs. The proofs that Lean checked are in
% supp-proofs.tex, each cut from this file unchanged (records/affine/split_q0448.py).
% `linkage check` reads every .tex file of the directory and matches a proof to its statement
% by the \ref of its heading, so a proof in the supplement is still compared with the blueprint.
% One rule decides the place of a proof: checked by Lean, the supplement; proved in print only,
% here. Two exceptions, both in print: the proofs of the classification (short, and the result
% the paper is about), and the proof of Proposition prop:profile-one-dimension-down (half a
% page, and the corollary proved in print reads it).
% A statement that changes status moves by one edit of the lists of section 1.1 and by moving
% its proof between this file and supp-proofs.tex.
\section{Notation, sources, and the proofs given in print}\label{app:isotropic-proofs}

This appendix has the notation (Table~\ref{tab:notation}), the sources of the statements
(\S\ref{sec:isotropic-sources}), the proofs of the classification
(\S\ref{app:classification-proofs}), five statements that the main text names and does not state
(\S\ref{app:further-statements}), and the proofs of the statements that are proved in print
only (\S\ref{app:print-proofs}). The proofs that Lean checked are in \suppl, except those of the
classification and that of Proposition~\ref{prop:profile-one-dimension-down}, which is in
\S\ref{app:print-proofs}.
% Q-0448, the blind draft review of the split (flag 4): the sentence that says how the numbers
% of the supplement read (\supplnumbers) is in "Plan of the paper", before the first S-number.

% Q-0428: the notation table of section 2.7, moved here as it stood (its rows for the new symbols were
% added in the same item).
% Q-0310: the notation table the brief asks for at the end of section 2 (review flag 25 and the clash
% map; records/affine/REVIEW-paper-whole-fixes-paper-iso.md). Paper-only; every row was checked
% against its first use. The second roles named are in shared statements and are owed at the
% blueprint (records/affine/GATE-paper-iso.md; Q-0311 added lambda_d, t and a, review flag 26).
% Q-0430: a longtable, as Table tab:in-print now is (sources.tex). As a float it could not stand
% on the appendix's first page and came out after that table, which is no float any more, so
% that Table 2 was printed before Table 1. Its rows and its caption are unchanged; N is the
% column p in the size of the table's text (main.tex).
{\renewcommand{\arraystretch}{0.88}
\begin{longtable}{@{}N{0.2\textwidth}N{0.76\textwidth}@{}}
\toprule
symbol & role \\
\midrule
\endhead
\bottomrule
\endfoot
\bottomrule
\addlinespace
\caption{Notation of \S\ref{sec:isotropic} and Appendix~\ref{app:isotropic-proofs}.}
\label{tab:notation}
\endlastfoot
$\xi$, $\omega$ & frequency on $\RR^d$, on the line \\
$s \le t$, $\chi$, $S_\lambda$ & scales; a gauge; the action of a dilation on the scales \\
$\lambda$, $\rho_\lambda$; $\lambda_d$ & a dilation factor; for $\lambda \in (0,1)$ in
  \ref{eq:self-decomposable}, the contraction ratio, and $\rho_\lambda$ the residual; the number
  $2^{-d/2}$ in Lemma~\ref{lem:threshold-monotone}, and the mean of $\log|\langle e,\eta\rangle|$
  over the sphere in Proposition~\ref{prop:isotropic-exponent-profile} (in C2, $\lambda$ is a
  measure on the sphere) \\
$\Phi_{s,t}$; $T_a$, $R_Q$, $D_\lambda$ & the operators of a family; on $L^1(\RR^d)$, the
  translation by $a$, the action of $Q \in \Orth(d)$ and the dilation by $\lambda$
  (Definition~\ref{def:isotropic-operator-family}; on the line $\tau_a$, $R$ and $D_\lambda$) \\
$\mu_{s,t}$; $\nu_{s,t}$, $\kappa_{s,t}$ & the kernels of a family on $\RR^d$; kernels on the line \\
$x$; $\mu^{\mathrm I}_{x,y}$ & a point of $\RR^d$ or of $\RR$; in $\mu^{\mathrm I}_{x,y}$, $x \le y$
  are scales of a time-causal family \\
$u_K$, $\Lambda$ & in C5, the uniform law on a convex set $K$, and the mixing measure \\
$F$, $F_d$, $F_{\mathrm I}$ & an exponent on the line \ref{eq:admissible-line}; the exponent of an
  isotropic origin datum on $\RR^d$; a causally admissible exponent \\
$a_{\mathrm V}$, $a$, $b_0$ & the Gaussian coefficient on the line; that of an isotropic origin datum,
  $a_{\mathrm V} = \tfrac12 a$ (in Lemma~\ref{lem:isotropic-exponent-log}(2), a dilation factor);
  the drift of a delay datum \\
$k$, $k_d$, $k_{\mathrm I}$ & the folded profile; the polar profile, the radial \Levy{} density times
  $|S^{d-1}|\,r^d$; the delay profile \\
$c$ & the profile at the origin, $k_d(0{+})$, and $k(0{+})$ on the line; in the uniqueness clause of
  Theorem~\ref{thm:isotropic-classification}, the ratio of two gauges \\
$\nu_d$, $f_d$; $p$, $q$ & the \Levy{} measure of an isotropic origin datum, and its radial
  density; the density of a kernel and its radial form, $p(x) = q(|x|)$ \\
$K$, $C_k$, $\varkappa_d$ & the slowly varying factor of the rate, the integral of the
  polar profile beyond radius $1$, and a constant of the dimension
  (Proposition~\ref{prop:isotropic-exponent-profile}); $K_\nu$ is a Bessel function, and in
  Lemma~\ref{lem:threshold-limit} $K$ is a number \\
$A$, $Q$, $B$; $R(d,c)$ & the constants of Proposition~\ref{prop:origin-rate-profile-d}, and that
  of Lemma~\ref{prop:origin-rate-d}. Elsewhere $A_\beta$ is the constant of
  Proposition~\ref{prop:origin-upper-bound-d}, $Q$ an orthogonal matrix
  (Definition~\ref{def:isotropic-operator-family}, Proposition~\ref{prop:radial-dictionary}), $B$
  a Borel set (Theorem~\ref{thm:origin-boundedness-d}(4) and C2), $R_Q$ the action of $Q$, and
  $R$ a random radius (Proposition~\ref{prop:radial-dictionary}) \\
$L$; $\mathcal{L}$ & the slowly varying factor of a transform, $\widehat\mu = |\xi|^{-c}L(|\xi|)$;
  the function $\int_s^1 K(y)\,dy/y$ of Proposition~\ref{prop:origin-rate-profile-d}(3) \\
$e$ & a unit vector of $\RR^d$, as in $\widehat\mu(e/s)$; in $e^{-C_k}$ and the like, the base of
  the exponential \\
$\Omega_d$, $\phi$, $X$ & in \S\ref{sec:isotropic-filtration}, the characteristic function of the
  uniform law on the unit sphere along a ray; that of a rotation-invariant random vector $X$
  (in \S\ref{sec:line-origin} $\phi_1$ is a density, and in
  Proposition~\ref{prop:isotropic-exponent-profile} $X$ is a real variable) \\
$\mathcal{A}_d$, $\mathcal{A}_\infty$, $\mathcal{S}$ & the classes of
  Proposition~\ref{prop:line-filtration}; their intersection; the subordinated slice \\
$b\,\Cin(\tau|\omega|)$; $\tau_a$ & a member of a $\Cin$ ray, of folded profile
  $b\,\mathbf{1}_{(0,\tau)}$; the translation by $a$ \\
$\gamma$, $\theta$; $\alpha$; $\vartheta$ & the \Matern{} index and range; the stable kernels, of
  index $2\alpha$; the shape of the Student-t kernels \\
$W$, $Z$; $T_1$ & a random delay and a standard normal vector; the scaled inverse-gamma delay \\
$w$, $w_d$, $M_d$, $m$ & the parameter of the separating family, its critical values,
  $m = (d+1)/2$ (Definition~\ref{def:filtration-thresholds}) \\
$\tau$ & the cut-off of a $\Cin$ ray and of a step datum; in \S\ref{sec:isotropic-sources},
  the time of a \Levy{} process, the kernel being its law at $\tau = 1$; a constant of
  Proposition~\ref{prop:origin-lower-bound-d} \\
$s$ & in Proposition~\ref{prop:origin-rate-profile-d} and in Lemma~\ref{prop:origin-rate-d} of
  \suppl, a distance from the centre of the kernel, not a scale \\
$g_t$, $g^{(d)}_u$ & the heat kernel at variance $t$ (here $t$ is not a scale); the centred
  Gaussian law of covariance $uI$ \\
\end{longtable}
}

% Q-0428, the fixes round: the one home of the attributions (sources.tex). It stands here and not
% at the end of section 3 because of the length.
% Sources --- the one home of the attributions of section 3 (Q-0428, the fixes round; the register
% review of sections 3.3 to 4, structural change 1, flags 1, 8, 16, 17, 18; the author's direction
% of 2026-10-12: attribution exact and made once). Paper-only: nothing here is shared with the
% blueprint. Until this round the table stood in the opening of section 3.3, before every
% statement it lists, with a third column that graded each statement ("restated", "extended",
% "proved"); the column is gone, and the cells are the source, its anchor and the case that the
% source covers. The statements of the body carry no sentence of attribution, but for two names a
% reader needs to use a statement (Watanabe after the shape; Halgreen for the Student-t clause).
%
% Anchors, each from the note that read it at the page (records/affine/), all confirmed by the
% citation audit of 2026-10-12 (CITATION-AUDIT-paper-iso-2026-10-12.md, 45 sites, all right):
% sato1999levy, Thm. 16.1, p. 99, by Sato 1991, p. 117, and the proof of Thm. 8.1(ii), p. 40
% (READING-sato1999levy-prior-art.md); hazod2001stable, Thm. 2.14.21, p. 383, with its limit
% condition (READING-hazod2001stable.md; the audit's note on row 38); knopova2013note, p. 126 and
% Lemma 2.6, p. 132 (READING-knopova2013note.md); grzywny2019asymptotic, Example 3, Thm. 3.5,
% Prop. 5.3, (5.4), by the numbers of the published article (READING-grzywny2019asymptotic.md);
% sato1978distribution, Lemma 5.1, p. 298, Thm. 1.7, p. 277 (READING-sato1978distribution.md);
% bingham1987regular, Thm. 3.7.3, pp. 162-163; watanabe1983isoperimetric, Proposition, p. 488
% (READING-watanabe1983isoperimetric.md: the law at each time of a symmetric Levy process whose
% Levy measure is isotropic unimodal; the audit's note on row 39); leonenko2013tauberian, Thms. 3
% and 4, arXiv v1 (READING-leonenko2013tauberian.md, with its addendum: the survey is where the
% theorems are collected; the citation is to the book); bingham1997abel, Thms. 1.1 and 1.2 and
% (1.4) of the preprint (READING-bingham1997abel.md, with its addendum: both for a function
% ultimately decreasing to zero); cygan2017asymptotic,
% Thms. 8 and 9, arXiv v2 (READING-cygan2017asymptotic.md); jurek1983integral, Thm. 3.2, p. 252,
% Thm. 4.1, p. 254, Thm. 7.2, p. 259 (READING-jurek1983integral-driving-process.md);
% sato1984operator, Thm. 4.1, p. 86; jurek1992operator, Cor. 1.7, pp. 1056-1057;
% jurek1993operator, Prop. 3.9.6, pp. 180-182; schoenberg1938metric, section 2.1 and Thm. 2,
% pp. 816-817; gneiting1999radial, p. 94; cambanis1981theory, Lemma 2, p. 371, Cor. 2, p. 372
% (the radius of a projection); fang1990symmetric, Thm. 2.10, p. 36 (the density of the
% projection itself, the audit's note on row 17; never its remark of p. 37); the cone paper,
% Lem. 4.12 and Rem. 4.10; trabs2014infinitely, Prop. 1, p. 58.
%
% Q-0430 (the primary sources, read at the page images; records/affine/
% CITATIONS-primary-sources-A.md and -B.md with their reading notes, and
% CITATION-CHANGES-q0430.md for every citation added, changed or left out, with the reason):
% leonenko1999limit, Thm. 1.4.3, p. 65 (the implication from the covariance function to the
% spectral function: 0 < alpha < n, n >= 1, no monotonicity), its proof on p. 104 from
% Lemma 1.4.11, p. 72; Thm. 1.4.4, pp. 66-67 (the book prints the monotonicity of the density
% as "decreasing for ||lambda|| >= k", the survey has (0, eps]: the cell says "a nonincreasing
% function of the radius" and quotes neither; nor the book's remark on c_5(1, alpha), p. 67,
% nor the constant of its Example 1.4.1, both misprinted); the credits of its Comments, p. 351;
% leonenko1991tauberian, Corollary of Lemma 2, p. 1542 of the translation (a pure power, the
% constant L(n) a_n(alpha)/Gamma(1 + alpha) with L(n) "a constant");
% soni1975slowlya, Thms. 1 and 2, p. 169, with the remark of p. 171; soni1975slowlyc,
% Thms. 17(a) and 18(a), pp. 613-614, the Hankel kernel on pp. 617-618 (its display of the
% coefficients c_j is off by a factor sqrt 2: no constant is taken from it). Bingham 1972 and
% Ridenhour and Soni 1974, to which Bingham and Inoue give their Theorem 1.2, are not cited:
% the first is a theorem for another transform, the second the converse implication under
% other conditions (READING-bingham1972tauberian.md, READING-ridenhour1974parseval.md).
% sato1983stationary, Thms. 5.1-5.3, pp. 546-547; jurek1982integral, Thms. 2.1 and 3.1,
% folios 388 and 390; wolfe1982continuous, Thm. 1, p. 306, Cor. 3, p. 307;
% daley2014dimension, Prop. 11 with (3.8), p. 1821, Lemma 12, p. 1822 (its Thm. 4, p. 1815, was
% read in the text layer only and is not cited).
% The versions cited (the author, 2026-10-12: the copies held are the ones the paper cites):
% cygan2017asymptotic and leonenko2013tauberian in the numbering of their arXiv versions (v2,
% v1), bingham1997abel in that of the Hokkaido preprint (no. 370), said once, in the caption.
% Their published numbering was not compared. grzywny2019asymptotic is cited by the statement
% numbers of the published article and never by a page.
%
% Q-0450 (the citation audit of 2026-10-13, CITATION-AUDIT-paper-iso-2026-10-13.md, 54 sites):
% two rows corrected. bogdan2014density, (19): p. 11 of the arXiv version held (v2), not p. 9,
% which has (17) and (18); read again at the page for this item. watanabe1983isoperimetric,
% Proposition, p. 488: its condition (3) is S = aI with a >= 0 and the Levy measure isotropic
% unimodal, so the row and the sentence after Proposition prop:isotropic-radial-monotone name
% the Gaussian covariance as well.

\subsection{Sources}\label{sec:isotropic-sources}

Table~\ref{tab:in-print} says where the main statements of \S\ref{sec:isotropic} are in print,
statement by statement, with the case that each source covers. Several of the sources are
written for a \Levy{} process (\S\ref{sec:line-bridge}): with a
law of exponent $F_d$ they study the laws with the transforms $e^{-\tau F_d}$, $\tau > 0$, the
laws of the process at the times $\tau$. The time is not the scale of this paper. The time
multiplies the exponent, and the scale rescales the frequency. A kernel is the law at time $1$,
and the law at time $\tau$ is a kernel with $\tau c$ in place of $c$.

% Q-0430: a longtable (it is longer than a page and breaks between rows). The struts of a row
% are those of the table's text size: \arraystretch is the ratio of the two line spacings.
{\renewcommand{\arraystretch}{0.81}
\begin{longtable}{@{}S{0.2\textwidth}S{0.76\textwidth}@{}}
\toprule
statement & where it is in print, and for which case \\
\midrule
\endhead
\bottomrule
\endfoot
\bottomrule
\addlinespace
\caption{Where the main statements of \S\ref{sec:isotropic} are in print, and for which case.
Three sources are cited in the numbering of the version that was read, which is not the
journal's: the Hokkaido preprint for \sscite{bingham1997abel}, and the arXiv versions for
\sscite{cygan2017asymptotic} and \sscite{leonenko2013tauberian}. The numbers of
\sscite{grzywny2019asymptotic} are those of the published article, and the pages of
\sscite{leonenko1991tauberian} those of the English translation.}
\label{tab:in-print}
\endlastfoot
Thm.~\ref{thm:isotropic-classification} and Cor.~\ref{cor:isotropic-scale-zero-kernels}, the
classification &
  Read on laws: a process with independent increments on $\RR^d$ that is self-similar up to a
  change of scale and a deterministic shift has self-decomposable laws, and every
  self-decomposable law but a point mass is the law at time $1$ of one
  (\sscite{sato1999levy}, Thm.~16.1, p.~99, due to Sato \sscite{sato1991selfsimilar}). On
  groups, a law is self-decomposable exactly when it is the limit of the kernels $\mu_{0,t}$ of
  a continuous hemigroup that a contracting one-parameter group of automorphisms maps into
  itself (\sscite{hazod2001stable}, Thm.~2.14.21, p.~383). \\
\addlinespace
Thm.~\ref{thm:origin-boundedness-d}, the threshold &
  On the line: \S\ref{sec:line-origin}. For the laws of a \Levy{} process on $\RR^d$: bounded
  when $c > d$, unbounded when $c \le d/2$ (\sscite{knopova2013note}, p.~126, and the proof of
  Lemma~2.6, p.~132). The threshold with its boundary case $c = d$, for two families: the
  \Matern{} kernels (the sources quoted in Example~\ref{ex:matern}, where two of the kernels
  above the threshold are given with their derivatives) and the laws with transform
  $(1 + |\xi|^\alpha)^{-\tau}$ (\sscite{grzywny2019asymptotic}, Example~3). \\
\addlinespace
Cor.~\ref{cor:origin-profile-exponent}, $c$ read off the exponent &
  Clause (1), for every infinitely divisible law: \sscite{sato1999levy}, proof of Thm.~8.1(ii),
  p.~40. Clause (2) for a finite $c$: on the line, an asymptotic form of the transform with its
  constant (\sscite{sato1978distribution}, Lemma~5.1, p.~298); in $\RR^d$ it follows from an
  equivalence for the \Levy{} processes whose laws are rotation invariant with radially
  nonincreasing densities (\sscite{grzywny2019asymptotic}, Thm.~3.5) by de Haan's theorem on
  functions whose increments are governed by a slowly varying function
  (\sscite{bingham1987regular}, Thm.~3.7.3, pp.~162--163). \\
\addlinespace
Prop.~\ref{prop:isotropic-radial-monotone}, the shape &
  For the law at each time of a symmetric \Levy{} process whose Gaussian covariance is a
  multiple of the identity and whose \Levy{} measure has a radially
  nonincreasing density: \sscite{watanabe1983isoperimetric}, Proposition, p.~488. That a
  self-decomposable law on $\RR^d$ which is not degenerate is absolutely continuous:
  \sscite{sato1999levy}, Thm.~27.13, p.~181, proved by Sato in 1982 (p.~196). \\
\addlinespace
Lemma~\ref{prop:origin-rate-d} of \suppl, the rate through the transform: clause (1) for
$c < d$, quoted before Proposition~\ref{prop:origin-rate-profile-d}, and clause (2) for $c = d$ &
  Clause (1), for a rotation-invariant probability law on $\RR^d$, $d \ge 1$, whose transform is
  $|\xi|^{-c}$ times a slowly varying factor, $0 < c < d$: the mass of a ball
  (\sscite{leonenko1999limit}, Thm.~1.4.3, p.~65, where the law is the spectral measure of a
  random field and the transform its covariance function), and the density when it is a
  nonincreasing function of the radius (the same book, Thm.~1.4.4, pp.~66--67). The book
  gives the two theorems for a transform asymptotic to a power to Leonenko and Olenko, and
  their proofs with a slowly varying factor to Olenko (p.~351). For a power the mass of a
  ball is \sscite{leonenko1991tauberian}, Corollary, p.~1542, with a constant that is not given
  in closed form. The survey \sscite{leonenko2013tauberian} collects both theorems, as its
  Thms.~3 and~4, with related ones. \\
\addlinespace
Lemma~\ref{prop:origin-rate-d}, through Hankel transforms &
  For $c > (d-3)/2$ the mass of a ball in clause (1) also follows from
  the Abel--Tauber theorem for Hankel transforms, in the direction from a function that is
  ultimately decreasing to zero to its transform. For a bounded function, which is the case
  $d \ge 3$: \sscite{soni1975slowlya}, Thm.~1, p.~169, with the remark on p.~171, and, as an
  equivalence for a function decreasing on the half-line, \sscite{soni1975slowlyc},
  Thm.~17(a), p.~613, with pp.~617--618 for the Hankel kernel. With local integrability in
  place of boundedness, and so for every $d$: \sscite{bingham1997abel}, Thm.~1.2. The mass of
  a ball in clause (2) follows from the boundary case of that theorem
  (\sscite{soni1975slowlya}, Thm.~2, p.~169, and \sscite{soni1975slowlyc}, Thm.~18(a),
  pp.~613--614, for a bounded function; \sscite{bingham1997abel}, Thm.~1.1 with (1.4)). \\
\addlinespace
Lemma~\ref{prop:origin-rate-d}, on the line, and the method of its proof &
  On the line: the theorem of Sato and Yamazato quoted in \S\ref{sec:line-origin}. The method
  of proof, an average of the transform against a positive kernel read as a Laplace transform of
  the mass of balls: with the Poisson kernel, for the spectral measure of a random field,
  \sscite{leonenko1999limit}, Lemma~1.4.11, p.~72, and the proof on p.~104; with the Gaussian
  kernel, in $\RR^d$, for the measure of the expected time that a \Levy{} process spends in a
  set, \sscite{cygan2017asymptotic}, Thm.~8 with its proof, and Thm.~9. The constant of the
  mass of a ball in clause (1), $R(d,c)\,|S^{d-1}|/c$, is the constant of these two theorems at
  the index $c$. \\
\addlinespace
Prop.~\ref{prop:origin-rate-profile-d}, the rate through the profile &
  On the line: the same theorem of Sato and Yamazato, with their lemma cited at
  Corollary~\ref{cor:origin-profile-exponent} above. In $\RR^d$, for the \Levy{} processes whose
  laws have radially nonincreasing densities and as the time $\tau$ tends to $0$: the density at
  $x$ is asymptotic to $\tau$ times the \Levy{} density at $x$ times the transform of the law at
  the frequencies of length $1/|x|$ (\sscite{grzywny2019asymptotic}, Thm.~3.9). The constant
  $R(d,\tau c)$ of the proposition for the law at time $\tau$ agrees with it in that limit. \\
\addlinespace
Props.~\ref{prop:origin-lower-bound-d} and~\ref{prop:origin-upper-bound-d}, the two bounds;
Lemma~\ref{lem:isotropic-exponent-log} &
  For the \Levy{} processes whose laws have radially nonincreasing densities, a lower bound and the
  comparison of a radially nonincreasing density with its average, on which the upper bound rests
  (\sscite{grzywny2019asymptotic}, Prop.~5.3 and (5.4)). The lemma on the line:
  \sscite{trabs2014infinitely}, Prop.~1, p.~58, and the same lemma of Sato and Yamazato. The
  two-sided estimates of \sscite{bogdan2014density} for such processes assume an exponent whose
  growth lies between two powers of positive order ((19), p.~11, of the arXiv version), which a
  finite $c$ excludes. \\
\addlinespace
Prop.~\ref{prop:isotropic-stages}, the stages &
  On the line, for the \Matern{} kernel of index $1$: the stage between two scales of ratio $r$
  as a mixture in which no displacement has the weight $r^2$ (\prog{B}{Prop.~7.2}, with (7.2)). \\
\addlinespace
Prop.~\ref{prop:radial-dictionary}, the dictionary &
  Clauses (1) and (2) are a case of the representation of a self-decomposable law as the limit
  law of the integrals $\int_0^t e^{-s}\,dY(s)$ of a \Levy{} process $Y$
  (\sscite{jurek1983integral}, Thms.~3.2 and~4.1, pp.~252 and~254; \sscite{sato1984operator},
  Thm.~4.1, p.~86): the case in which $Y$ moves by isolated jumps, at the rate $c$ and with the
  law of $X$. The representation is also in \sscite{sato1983stationary}, Thms.~5.1--5.3,
  pp.~546--547 (the law of a stationary process of Ornstein--Uhlenbeck type), in
  \sscite{jurek1982integral}, Thms.~2.1 and~3.1, pp.~388 and~390 (operator-self-decomposable
  laws on a Banach space), and, on the line, in \sscite{wolfe1982continuous}, Thm.~1, p.~306,
  and Cor.~3, p.~307 (the limit in law of $e^{-\gamma t}\int_0^t e^{\gamma u}\,dB(u)$,
  $\gamma > 0$). In radial form: \sscite{jurek1992operator}, Cor.~1.7, pp.~1056--1057, restated as
  \sscite{jurek1993operator}, Prop.~3.9.6, pp.~180--182. The classes of the radial
  characteristic functions on $\RR^m$ decrease in $m$, with the mixtures of Gaussians as their
  intersection (\sscite{schoenberg1938metric}, \S2.1 and Thm.~2, pp.~816--817), and strictly
  (\sscite{gneiting1999radial}, p.~94). \\
\addlinespace
Prop.~\ref{prop:profile-one-dimension-down} and Cor.~\ref{cor:dictionary-filtration}, the test
on the \Levy{} measure &
  For probability laws, the radius of a projection of a rotation-invariant law
  (\sscite{cambanis1981theory}, Lemma~2, p.~371, and Cor.~2, p.~372) and the density of the
  projection (\sscite{fang1990symmetric}, Thm.~2.10, p.~36). See also
  \sscite{daley2014dimension}: Prop.~11 with (3.8), p.~1821, for the radius law of a radial
  characteristic function in a lower dimension, the law of its radius times an independent
  variable with values in $(0,1)$. The formula of Lemma~\ref{lem:shadow-polar-profile} is (3.8)
  with the \Levy{} measure in place of the law. And Lemma~12, p.~1822: $\Omega_d$ is a
  characteristic function on $\RR^d$ and not on $\RR^{d+1}$, which is the last assertion of
  Corollary~\ref{cor:dictionary-filtration}(2). The representation of a
  self-decomposable law of which the $\Cin$ rays are the elementary terms on the line:
  \sscite{jurek1983integral}, Thm.~7.2, p.~259. \\
\addlinespace
Thm.~\ref{thm:separating-family-critical-values}, the separating family &
  The family, the members $w = 8$ and $w = 4$ as witnesses in dimensions $1$ and $2$, and the
  first critical values, computed numerically: \prog{B}{Lem.~4.12 and Rem.~4.10}. \\
\end{longtable}
}


\subsection{The proofs of the classification}\label{app:classification-proofs}

The necessity half of Theorem~\ref{thm:isotropic-classification} is proved in
\S\ref{sec:isotropic-reduction}. This subsection has the other proofs of the classification:
those of Lemma~\ref{lem:radial-marginal}, of the sufficiency and uniqueness halves of the
theorem, of Corollary~\ref{cor:isotropic-classification-operators}, of
Lemma~\ref{lem:isotropic-realisation}, of Proposition~\ref{prop:isotropic-polar-reading} and of
Corollary~\ref{cor:isotropic-scale-zero-kernels}.
The proofs of Lemmas~\ref{lem:operators-to-kernels} and~\ref{lem:isotropic-marginal-family},
which carry arguments of the line paper from the line to $\RR^d$, are in \suppl.

\begin{proof}[Proof of Lemma~\ref{lem:radial-marginal}]
\emph{(1).} $\pi_*\mu$ is a probability measure because $\pi$ is measurable, and for $\omega \in \RR$
\[
  \widehat\kappa(\omega) = \int_{\RR} e^{i\omega y}\,\kappa(dy)
   = \int e^{i\omega\langle e, x\rangle}\,\mu(dx) = \widehat\mu(\omega e) .
\]
A rotation-invariant law has a radial transform: if $|\eta| = |\eta'|$, let $Q$ be the identity when
$\eta = \eta'$ and otherwise the reflection in the hyperplane orthogonal to $\eta - \eta'$, an
element of $\Orth(d)$ with $Q^{\tp} = Q$ and $Q\eta = \eta'$; then
$\widehat\mu(\eta) = \widehat{Q_*\mu}(\eta) = \widehat\mu(Q^{\tp}\eta) = \widehat\mu(\eta')$.
Since $|\,|\xi|e\,| = |\xi|$, this gives $\widehat\mu(\xi) = \widehat\mu(|\xi|e) =
\widehat\kappa(|\xi|)$ for every $\xi \in \RR^d$. Since $-I \in \Orth(d)$ and $\pi(-x) = -\pi(x)$,
the image of $\kappa$ under $\omega \mapsto -\omega$ is $\pi_*((-I)_*\mu) = \kappa$, so $\kappa$ is
symmetric and $\widehat\kappa$ is real; hence so is $\widehat\mu$. For another unit vector $e'$ with
marginal $\kappa'$, $|\omega e'| = |\omega e|$ gives $\widehat{\kappa'}(\omega) = \widehat\mu(\omega e')
= \widehat\mu(\omega e) = \widehat\kappa(\omega)$ for every $\omega \in \RR$, so $\kappa' = \kappa$
by uniqueness of the transform on $\RR$.

\emph{(2).} If $\mu = \delta_0$ then $\pi_*\mu = \delta_0$. Conversely, if $\kappa = \delta_0$ then
$\widehat\kappa \equiv 1$, so $\widehat\mu \equiv 1 = \widehat{\delta_0}$ by (1), and
$\mu = \delta_0$ by uniqueness of the transform on $\RR^d$.

\emph{(3).} $\mu * \mu'$ is the image of $\mu \otimes \mu'$ under $(x,y) \mapsto x + y$, and
$\pi(x+y) = \pi(x) + \pi(y)$, so $\pi_*(\mu * \mu')$ is the image of $\pi_*\mu \otimes \pi_*\mu'$
under addition on $\RR$, which is $\pi_*\mu * \pi_*\mu'$. The other two identities are
$\pi(\lambda x) = \lambda\pi(x)$ and $\pi(0) = 0$.
\end{proof}

\begin{proof}[Proof of Theorem~\ref{thm:isotropic-classification}, sufficiency and uniqueness]
\emph{Sufficiency.} Let $\chi$ be a gauge and $F \in \mathcal{A}_d$ with $F \not\equiv 0$, and let
$(\mu_{s,t})$ be a family of probability measures on $\RR^d$ satisfying
\ref{eq:isotropic-classification}. First, $F$ is
continuous on $[0,\infty)$ with $F(0) = 0$: on $0 \le r \le \Omega$ the integrand of
\ref{eq:admissible-line} is at most $\max\{2,\Omega^2/2\}\,(1 \wedge x^2)\,k(x)/x$, which is
integrable by the two conditions on $k$, so dominated convergence applies. The axioms are read off
the transform, a probability measure on $\RR^d$ being determined by its
transform. (I1): the exponents of $\mu_{s,t}$ and $\mu_{r,s}$ add up to that of $\mu_{r,t}$, and
the exponent of $\mu_{t,t}$ is $0$. (I2): $F$ is continuous, and so is $\chi$, an increasing
bijection of $[0,\infty)$. (I3): $\widehat{Q_*\mu_{s,t}}(\xi) = \widehat\mu_{s,t}(Q^{\tp}\xi)$ and
$|Q^{\tp}\xi| = |\xi|$. (I4): with $S_\lambda := \chi^{-1}(\lambda\,\chi(\cdot))$, an increasing
bijection of $[0,\infty)$,
$\widehat{(\lambda I)_*\mu_{s,t}}(\xi) = \widehat\mu_{s,t}(\lambda\xi)
= \widehat\mu_{S_\lambda s,S_\lambda t}(\xi)$. (ND): let $s < t$ and suppose
$\mu_{s,t} = \delta_0$, so that $F(\chi(t)r) = F(\chi(s)r)$ for every $r \ge 0$. If $s = 0$ this
says $F(\chi(t)r) = 0$ for every $r$, and $\chi(t) > 0$ gives $F \equiv 0$. If $s > 0$ it says
$F(r) = F(\varkappa r)$ for every $r$, with $\varkappa := \chi(s)/\chi(t) \in (0,1)$. Iterating,
$F(r) = F(\varkappa^nr) \to F(0) = 0$, so again $F \equiv 0$ (\restated{V}{Lem.~6.3}{sec:line-facts}). Both contradict
$F \not\equiv 0$.

\emph{Uniqueness.} Let $(\chi',F')$ also satisfy \ref{eq:isotropic-classification}. Both exponents
are real, so taking logarithms of the kernels from scale zero, $F(\chi(t)r) = F'(\chi'(t)r)$ for
all $t, r \ge 0$, since $\chi(0) = \chi'(0) = 0$ and $F(0) = F'(0) = 0$. At $t' := \chi'^{-1}(1)$ this reads $F'(r) = F(c^{-1}r)$ with
$c^{-1} := \chi(t') > 0$. Hence $F(\chi(t)r) = F(c^{-1}\chi'(t)r)$ for all $t$ and $r$. If
$\chi(t) \ne c^{-1}\chi'(t)$ for some $t > 0$, then $F(r) = F(\varkappa r)$ for every $r$ with
$\varkappa := c^{-1}\chi'(t)/\chi(t) \ne 1$, and $F \equiv 0$ as in the previous paragraph (with
$\varkappa^{-1}$ in place of $\varkappa$ if $\varkappa > 1$), a contradiction. So
$\chi' = c\,\chi$. When $\chi(1) = \chi'(1) = 1$ this gives $c = 1$.
\end{proof}

% Q-0414, round 2: the proof of Corollary cor:isotropic-classification-operators, moved here
% from section 3.2 (records/affine/edits_q0414_round2_moves.py); the body is unchanged.
\begin{proof}[Proof of Corollary~\ref{cor:isotropic-classification-operators}]
\emph{Necessity.} Let $(\Phi_{s,t})$ be an isotropic cascade measurement family in dimension $d$.
By Lemma~\ref{lem:operators-to-kernels}(1) each $\Phi_{s,t}$ is convolution by a probability measure
$\mu_{s,t}$ on $\RR^d$, and by Lemma~\ref{lem:operators-to-kernels}(3) the family $(\mu_{s,t})$ is an
isotropic cascade family in dimension $d$. The necessity direction of
Theorem~\ref{thm:isotropic-classification} gives a gauge $\chi$ and an $F \in \mathcal{A}_d$ with
$F \not\equiv 0$ satisfying \ref{eq:isotropic-classification}.

\emph{Sufficiency.} Let $\chi$, $F$ and $(\mu_{s,t})$ be as in the statement. By the sufficiency
direction of Theorem~\ref{thm:isotropic-classification} the family $(\mu_{s,t})$ is an isotropic cascade
family in dimension $d$. By Lemma~\ref{lem:operators-to-kernels}(4) the operators
$f \mapsto \mu_{s,t} * f$, which are the $\Phi_{s,t}$, form an isotropic cascade measurement family
in dimension $d$.

\emph{Uniqueness.} Let $\chi'$, $F'$ and $(\mu'_{s,t})$ be as in the last sentence of the
statement. For $0 \le s \le t$ the probability measures $\mu_{s,t}$ and $\mu'_{s,t}$ give the same
operator $\Phi_{s,t}$, so they are equal, by the uniqueness clause of
Lemma~\ref{lem:operators-to-kernels}(1). Hence $(\chi', F')$ satisfies
\ref{eq:isotropic-classification} for the family $(\mu_{s,t})$, and the uniqueness clause of
Theorem~\ref{thm:isotropic-classification} gives $c$.
\end{proof}

\begin{proof}[Proof of Lemma~\ref{lem:isotropic-realisation}]
Let $\mu$ be a self-decomposable probability measure on $\RR^d$ with
$\widehat\mu(\xi) = e^{-F(|\xi|)}$, which the definition of $\mathcal{A}_d$ provides. For
$c \in (0,1)$ let $\rho_c$ be a residual of $\mu$. As $\widehat\mu(c\,\xi) = e^{-F(c|\xi|)} \ne 0$,
\ref{eq:self-decomposable} forces $\widehat\rho_c(\xi) = e^{-(F(|\xi|) - F(c|\xi|))}$; put
$\rho_0 := \mu$ and $\rho_1 := \delta_0$, so that this formula holds for every $c \in [0,1]$, as
$F(0) = 0$. For $0 \le s \le t$ put
\[
  \mu_{s,t} := \bigl(\chi(t)I\bigr)_*\,\rho_{\chi(s)/\chi(t)} ,
\]
with $\chi(s)/\chi(t) := 0$ when $\chi(t) = 0$, that is, when $s = t = 0$; then $\mu_{0,0}$ is
the image of $\mu$ under the zero map, $\delta_0$. Since
$0 \le \chi(s) \le \chi(t)$, the index lies in $[0,1]$, and
\[
  \widehat\mu_{s,t}(\xi) = \widehat\rho_{\chi(s)/\chi(t)}\bigl(\chi(t)\xi\bigr)
  = \exp\bigl[-\bigl(F(\chi(t)|\xi|) - F(\chi(s)|\xi|)\bigr)\bigr] ,
\]
because $\chi(s)/\chi(t)\cdot\chi(t)|\xi| = \chi(s)|\xi|$, both sides being $0$ when
$\chi(t) = 0$. Each $\mu_{s,t}$ is a probability measure, the image of one, so by the sufficiency
direction of Theorem~\ref{thm:isotropic-classification} $(\mu_{s,t})$ is an isotropic cascade family, and
it satisfies \ref{eq:isotropic-classification} with $(\chi, F)$ by construction.
\end{proof}

% Q-0414, round 2: Lemma lem:polar-decomposition, moved here from section 3.2 with its proof
% (records/affine/edits_q0414_round2_moves.py). The lead sentence is paper prose.
The proofs of this paper integrate against the \Levy{} measure of an isotropic origin datum in
polar coordinates. The identity \ref{eq:polar-decomposition} of
Definition~\ref{def:isotropic-origin-datum} is stated once, as a lemma.

% shared with blueprint lem:polar-decomposition
\begin{lemma}[the polar decomposition of an isotropic origin datum's \Levy{} measure]\label{lem:polar-decomposition}
Let $d \ge 1$ and let $(a, k_d)$ be an isotropic origin datum in dimension $d$ with \Levy{} measure
$\nu_d$. Then \ref{eq:polar-decomposition} holds for every nonnegative Borel $g$ on $\RR^d$
vanishing at the origin, $\sigma_{d-1}$ being the surface measure of $S^{d-1}$ and $|S^{d-1}|$ its
surface area.
\end{lemma}

\begin{proof}
By definition $\nu_d$ has the density $f_d(|x|) = k_d(|x|)(|S^{d-1}||x|^d)^{-1}$ on
$\RR^d \setminus \{0\}$. In polar coordinates $x = r\eta$, $dx = r^{d-1}dr\,\sigma_{d-1}(d\eta)$, so
for $g \ge 0$ Borel and vanishing at the origin
$\int g\,d\nu_d = \int_{S^{d-1}}\sigma_{d-1}(d\eta)\int_0^\infty g(r\eta)\,
k_d(r)(|S^{d-1}|r^d)^{-1}r^{d-1}\,dr$, which is \ref{eq:polar-decomposition}.
\end{proof}

\begin{proof}[Proof of Proposition~\ref{prop:isotropic-polar-reading}]
\emph{(1), forward.} Let $\mu$ be self-decomposable on $\RR^d$ with
$\widehat\mu(\xi) = e^{-F(|\xi|)}$. By \ledger{A4} it is infinitely divisible, and by \ledger{A3}(i)
it has a generating triplet $(A,\nu,\gamma)$. For $Q \in \Orth(d)$,
$\widehat{Q_*\mu}(\xi) = \widehat\mu(Q^{\tp}\xi) = \widehat\mu(\xi)$, so $Q_*\mu = \mu$. The
unit ball is $Q$-invariant, so substituting $x = Q^{\tp}y$ in the representation shows that
$Q_*\mu$ has triplet $(QAQ^{\tp}, Q_*\nu, Q\gamma)$. By uniqueness, \ledger{A3}(ii),
$QAQ^{\tp} = A$ and $Q_*\nu = \nu$ for every $Q$. For unit vectors $e, e'$ choose $Q$ with
$Qe = e'$. Then $\langle e',Ae'\rangle = \langle e, Q^{\tp}AQ\,e\rangle = \langle e,Ae\rangle$, so
$\langle\xi,A\xi\rangle = a|\xi|^2$ with $a := \langle e,Ae\rangle \ge 0$. The real part of
$e^{i\theta} - 1 - i\theta\mathbf 1_D$ is $\cos\theta - 1$, so taking moduli in the representation,
\[
  F(|\xi|) = -\log\bigl|\widehat\mu(\xi)\bigr|
   = \tfrac12 a|\xi|^2 + \int_{\RR^d}\bigl(1 - \cos\langle\xi,x\rangle\bigr)\,\nu(dx) .
\]
Since $\mu$ is self-decomposable, \ledger{A2} gives $\nu(B) = \int_{S}\lambda(d\eta)\int_0^\infty
\mathbf 1_B(r\eta)\,k_\eta(r)\,dr/r$ with $\lambda$ a finite measure on $S = S^{d-1}$ and
$k_\eta(r)$ measurable in $\eta$ and nonincreasing in $r$.
A nonincreasing $k_\eta$ differs from its right limit
$k_\eta(r+) = \sup\{k_\eta(q) : q \in \mathbb Q,\ q > r\}$ at countably many $r$ only, a
Lebesgue-null set. Replacing $k_\eta(r)$ by $k_\eta(r+)$ therefore changes no integral in the
display, and makes the kernel jointly measurable in $(\eta,r)$, as a countable supremum of
measurable functions. Taking $B = \{x : |x| \in C\}$ for Borel $C \subset (0,\infty)$ and
exchanging the integrals, the image of $\nu$ under $x \mapsto |x|$ is
\[
  \nu\bigl(\{x : |x| \in C\}\bigr) = \int_C \bar k(r)\,\frac{dr}{r} ,
  \qquad \bar k(r) := \int_S k_\eta(r+)\,\lambda(d\eta) ,
\]
and $\bar k$ is nonincreasing with values in $[0,\infty]$. Integrating $|x|^2 \wedge 1$ against
$\nu$ through this image gives
$\int_0^\infty(r^2 \wedge 1)\bar k(r)\,dr/r = \int(|x|^2 \wedge 1)\,\nu(dx) < \infty$, which
contains the pair of conditions of Definition~\ref{def:isotropic-origin-datum}. They force $\bar k < \infty$
on $(0,\infty)$, because a nonincreasing function infinite at some $r_0$ is infinite on $(0,r_0]$.
So $(a, \bar k)$ is an isotropic origin datum. Let $\nu_d$ be its \Levy{} measure, the right side
of \ref{eq:polar-decomposition}. By Lemma~\ref{lem:polar-decomposition} the image of $\nu_d$ under
$x \mapsto |x|$ is $\bar k(r)\,dr/r$ on $(0,\infty)$, the same as that of $\nu$.

It remains to see that $\nu_d = \nu$. Both are invariant under every $Q \in \Orth(d)$ ($\nu$ by the
first paragraph, $\nu_d$ by the invariance of $\sigma_{d-1}$), both vanish at the origin, and both
are finite on each set $U_\varepsilon = \{x : |x| > \varepsilon\}$, $\varepsilon > 0$. A finite
measure $m$ invariant under $\Orth(d)$ has a transform that depends only on its radial image:
$\widehat m(Q\xi) = \widehat m(\xi)$ for every $Q$, a reflection carries $\xi$ to $|\xi|\theta$ for
any unit $\theta$, and averaging over $\theta$ against $\sigma_{d-1}/|S^{d-1}|$ and exchanging the
integrals gives
\[
  \widehat m(\xi) = \frac{1}{|S^{d-1}|}\int_{[0,\infty)}\Phi(|\xi|,r)\,(|\cdot|_*m)(dr) ,
  \qquad
  \Phi(s,r) := \int_S e^{i s r\langle e_1,\theta\rangle}\,\sigma_{d-1}(d\theta) ,
\]
the inner sphere integral $\int_S e^{i s\langle x,\theta\rangle}\sigma_{d-1}(d\theta)$ depending on
$x$ only through $|x|$, again by the invariance of $\sigma_{d-1}$. Applied to the restrictions of
$\nu$ and $\nu_d$ to $U_\varepsilon$, which are invariant and have the same radial image, this gives
equal transforms, so the restrictions are equal. Letting $\varepsilon \downarrow 0$ along
$1/(n+1)$ gives $\nu = \nu_d$ off the origin, hence everywhere. The display for $F$ is therefore
$F(|\xi|) = \tfrac12 a|\xi|^2 + \int(1 - \cos\langle\xi,x\rangle)\,\nu_d(dx)$, the exponent of the
datum $(a, \bar k)$.

\emph{(1), converse.} Let $(a,k_d)$ be a datum with $F_d = F(|\cdot|)$. Its \Levy{} measure
$\nu_d$ satisfies $\int(|x|^2 \wedge 1)\,d\nu_d < \infty$ by the two conditions, so by
\ledger{A3}(iii) there is an infinitely divisible $\mu$ with triplet $(aI, \nu_d, 0)$. As $\nu_d$
is symmetric, the odd part of the integrand integrates to $0$ and
$\widehat\mu = e^{-F_d}$: $\mu$ is the datum's law. Its \Levy{} measure has the form of
\ledger{A2} with $\lambda = \sigma_{d-1}$ and $k_\eta = k_d/|S^{d-1}|$ for every $\eta$,
nonincreasing in $r$. So $\mu$ is self-decomposable by \ledger{A2}.

\emph{(1), uniqueness.} Two data with the same exponent have the same law, hence by \ledger{A3}(ii)
the same Gaussian matrix $aI$ and the same \Levy{} measure. The images of the two \Levy{}
measures under $x \mapsto |x|$ are then equal; by Lemma~\ref{lem:polar-decomposition} they are
$k_d(r)\,dr/r$ and $k_d'(r)\,dr/r$ on $(0,\infty)$, so $k_d = k_d'$ almost everywhere.

\emph{(2).} Let $\mu$ be the law of clause (1), $e_1$ the first basis vector and $\pi(x) :=
\langle e_1,x\rangle$. Then $\pi_*\mu$ is self-decomposable on $\RR$. If $\rho_\varkappa$ is a
residual of $\mu$ for the ratio $\varkappa \in (0,1)$, then $\pi_*\rho_\varkappa$ is one of
$\pi_*\mu$:
\[
  \widehat{\pi_*\mu}(\omega) = \widehat\mu(\omega e_1)
   = \widehat\mu(\varkappa\omega e_1)\,\widehat\rho_\varkappa(\omega e_1)
   = \widehat{\pi_*\mu}(\varkappa\omega)\,\widehat{\pi_*\rho_\varkappa}(\omega) .
\]
Its exponent is $F(|\omega|)$.
Clause (1) at $d = 1$ gives a datum $(a',k_1)$ in dimension $1$ with, as $|S^0| = 2$ and the
integrand is even,
\[
  F(|\omega|) = \tfrac12 a'\omega^2 + \int_{\RR}\bigl(1 - \cos\omega x\bigr)\frac{k_1(|x|)}{2|x|}\,dx
   = \tfrac12 a'\omega^2 + \int_0^\infty\bigl(1 - \cos\omega x\bigr)\,k_1(x)\,\frac{dx}{x} ,
\]
which is \ref{eq:admissible-line} with $a_{\mathrm V} = \tfrac12 a'$ and $k = k_1$. The two
integrability conditions are the same pair. With the definition of $\mathcal{A}_d$ in
Proposition~\ref{prop:line-filtration} this is the description of the class.

\emph{(3).} By Theorem~\ref{thm:isotropic-classification} and clause (2), the family is isotropic exactly
when \ref{eq:isotropic-classification} holds for a gauge and an $F \not\equiv 0$ with $F(|\cdot|)$
the exponent $F_d$ of a datum. The displayed $F(r)$ is $F_d(re_1)$, and $F_d$ is radial. Finally
$F \equiv 0$ exactly when the datum is trivial. If $a = 0$ and $k_d = 0$ almost everywhere then
$F_d \equiv 0$. Conversely $F_d \equiv 0$ makes the law $\delta_0$, whose triplet is $(0,0,0)$,
so $a = 0$ and $k_d = 0$ almost everywhere by the uniqueness of clause (1).
\end{proof}

% Q-0448: the proof of record of cor:isotropic-scale-zero-kernels (Q-0434), written by
% records/affine/transcribe-q0448-nodes.py and pinned. It stands after the proof of
% Proposition prop:isotropic-polar-reading, which it cites.
\begin{proof}[Proof of Corollary~\ref{cor:isotropic-scale-zero-kernels}]
\emph{A kernel from scale zero satisfies (S).} Let $(\mu_{s,t})$ be an isotropic cascade family in
dimension $d$ and let $t > 0$. Then $\mu_{0,t}$ is rotation invariant by (I3), and
$\mu_{0,t} \ne \delta_0$ by (ND). By Theorem~\ref{thm:isotropic-classification} there are a gauge $\chi$
and an $F \in \mathcal{A}_d$ with \ref{eq:isotropic-classification}, so that
$\widehat\mu_{0,t}(\xi) = e^{-F(\chi(t)|\xi|)}$. By the definition of $\mathcal{A}_d$
(Proposition~\ref{prop:line-filtration}) there is a self-decomposable probability measure $\mu'$ on
$\RR^d$ with $\widehat{\mu'}(\xi) = e^{-F(|\xi|)}$. Put $\lambda := \chi(t) > 0$. Then
$\widehat{(\lambda I)_*\mu'}(\xi) = \widehat{\mu'}(\lambda\xi) = e^{-F(\lambda|\xi|)}
= \widehat\mu_{0,t}(\xi)$, so $\mu_{0,t} = (\lambda I)_*\mu'$, a probability measure on $\RR^d$
being determined by its transform. The image of a self-decomposable law under $\lambda I$ is
self-decomposable: if $\rho_c$ is a residual of $\mu'$ for the ratio $c \in (0,1)$, then
$\widehat{\mu'}(\lambda\xi) = \widehat{\mu'}(c\lambda\xi)\,\widehat\rho_c(\lambda\xi)$ is
\ref{eq:self-decomposable} for $(\lambda I)_*\mu'$ with the residual $(\lambda I)_*\rho_c$. So
$\mu_{0,t}$ satisfies (S).

\emph{The exponent of a law with (S).} Let $\mu$ satisfy (S). By Lemma~\ref{lem:radial-marginal}(1)
its transform is real and radial. It has no zero: $\mu$ is infinitely divisible by \ledger{A4},
so by \ledger{A3}(i) it has a generating triplet, and its transform is the exponential of a
complex-valued function. A real continuous function on $\RR^d$ that has no zero and takes the
value $1$ at the origin is positive. Fix a unit vector $e$ and put
$F(r) := -\log\widehat\mu(re)$ for $r \ge 0$. Then $F(0) = 0$, $F \ge 0$ because
$|\widehat\mu| \le 1$, and $\widehat\mu(\xi) = \widehat\mu(|\xi|e) = e^{-F(|\xi|)}$ for every
$\xi$, the transform being radial. So $\xi \mapsto F(|\xi|)$ is the exponent of the
self-decomposable law $\mu$, and $F$ is admissible by
Proposition~\ref{prop:isotropic-polar-reading}(2). Hence $F \in \mathcal{A}_d$. If $F$ vanished on
$[0,\infty)$ the transform of $\mu$ would be $1$ and $\mu$ would be $\delta_0$, so
$F \not\equiv 0$. A function $F'$ with $\widehat\mu(\xi) = e^{-F'(|\xi|)}$ for every $\xi$ has
$F'(r) = -\log\widehat\mu(re) = F(r)$ for every $r \ge 0$.

\emph{Clause (3), the equivalence.} Let $\mu$ satisfy (S), let $F$ be its exponent, let
$t_1 > 0$, and let $(\mu_{s,t})$ be a family of probability measures on $\RR^d$. Suppose first
that it satisfies \ref{eq:isotropic-classification} with $F$ and a gauge $\chi$ with
$\chi(t_1) = 1$. By the sufficiency direction of Theorem~\ref{thm:isotropic-classification}, $F$ being in
$\mathcal{A}_d$ and not identically $0$, it is an
isotropic cascade family, and
$\widehat\mu_{0,t_1}(\xi) = e^{-F(\chi(t_1)|\xi|)} = e^{-F(|\xi|)} = \widehat\mu(\xi)$, because
$\chi(0) = 0$ and $F(0) = 0$. So $\mu_{0,t_1} = \mu$. Conversely let $(\mu_{s,t})$ be an
isotropic cascade family with $\mu_{0,t_1} = \mu$. By the necessity direction of
Theorem~\ref{thm:isotropic-classification} there are a gauge $\chi'$ and an $F' \in \mathcal{A}_d$ with
\ref{eq:isotropic-classification}. Put $c := \chi'(t_1) > 0$. The kernel from scale zero to
scale $t_1$ gives $e^{-F'(c|\xi|)} = \widehat\mu(\xi) = e^{-F(|\xi|)}$, that is $F'(cr) = F(r)$
for every $r \ge 0$. Then $\chi := c^{-1}\chi'$ is a gauge with $\chi(t_1) = 1$, and
$F'(\chi'(t)|\xi|) = F'(c\,\chi(t)|\xi|) = F(\chi(t)|\xi|)$ for every $t \ge 0$. So the family
satisfies \ref{eq:isotropic-classification} with $(\chi, F)$.

\emph{Clause (3), the rest.} Let two gauges $\chi$ and $\chi''$ with
$\chi(t_1) = \chi''(t_1) = 1$ both satisfy \ref{eq:isotropic-classification} with $F$ for one
family. By the uniqueness clause of Theorem~\ref{thm:isotropic-classification}, applied with $F' = F$,
which vanishes at $0$, there is $c > 0$ with $\chi'' = c\,\chi$, and at $t_1$ this reads $c = 1$.
Let $\chi$ be any gauge with $\chi(t_1) = 1$. By Lemma~\ref{lem:isotropic-realisation} there is a
family of probability measures that satisfies \ref{eq:isotropic-classification} with
$(\chi, F)$, and by the equivalence it is an isotropic cascade family with $\mu_{0,t_1} = \mu$.
It is the only one, because \ref{eq:isotropic-classification} gives the transform of every
$\mu_{s,t}$. For $t > 0$,
$\widehat{(\chi(t)I)_*\mu}(\xi) = \widehat\mu(\chi(t)\xi) = e^{-F(\chi(t)|\xi|)}
= \widehat\mu_{0,t}(\xi)$, so $\mu_{0,t} = (\chi(t)I)_*\mu$.

\emph{Clause (1).} One direction is the first paragraph. For the other let $\mu$ satisfy (S).
Clause (3) with $t_1 = 1$ and the gauge $\chi = \mathrm{id}$ gives an isotropic cascade family
with $\mu_{0,1} = \mu$.

\emph{Clause (2).} Let $(\Phi_{s,t})$ be an isotropic cascade measurement family in dimension
$d$ and let $t > 0$ with $\Phi_{0,t}f = \mu * f$ for every $f$. By
Lemma~\ref{lem:operators-to-kernels}(1) each $\Phi_{s,t}$ is convolution by a probability measure
$\mu_{s,t}$, which it determines, and by Lemma~\ref{lem:operators-to-kernels}(3) the family
$(\mu_{s,t})$ is an isotropic cascade family in dimension $d$. As $\mu_{0,t}$ and $\mu$ give the
same operator, $\mu_{0,t} = \mu$, and (S) holds by clause (1). Conversely let $\mu$ satisfy (S).
By clause (1) there are an isotropic cascade family $(\mu_{s,t})$ and a $t > 0$ with
$\mu_{0,t} = \mu$. By Lemma~\ref{lem:operators-to-kernels}(4) the operators
$\Phi_{s,t}f := \mu_{s,t} * f$ form an isotropic cascade measurement family in dimension $d$,
and $\Phi_{0,t}f = \mu * f$.

\emph{Clause (4).} Let the two families be as in the clause and put $\mu := \mu_{0,t_1}$. By
the first paragraph $\mu$ satisfies (S). Let $F$ be its exponent, and let $\chi$ and $\chi'$ be
the gauges of the two families by clause (3), so that $\chi(t_1) = \chi'(t_1) = 1$. Then
$\theta := \chi^{-1} \circ \chi'$ is an increasing bijection of $[0,\infty)$ with
$\theta(t_1) = t_1$ and $\chi \circ \theta = \chi'$. Let $0 \le s \le t$, so that
$0 \le \theta(s) \le \theta(t)$. By \ref{eq:isotropic-classification} the transforms of
$\mu'_{s,t}$ and of $\mu_{\theta(s),\theta(t)}$ are both
$\exp[-(F(\chi'(t)|\xi|) - F(\chi'(s)|\xi|))]$, so $\mu'_{s,t} = \mu_{\theta(s),\theta(t)}$.
\end{proof}

\subsection{Statements named in the main text}\label{app:further-statements}

The five statements of this subsection are named in \S\ref{sec:isotropic} or used in the proofs
of \S\ref{app:print-proofs}. Lean checked them, and their proofs are in \suppl. The first
bounds the transform of a kernel with a finite $c$ from below.

% shared with blueprint lem:isotropic-exponent-log
\begin{lemma}[a profile bounded at the origin gives a logarithmic exponent]\label{lem:isotropic-exponent-log}
Let $d \ge 1$, let $(0,k_d)$ be an isotropic origin datum in dimension $d$ with Gaussian
coefficient $0$ and $c = k_d(0{+}) < \infty$, let $F_d$ be its exponent, and put
$C_k := \int_1^\infty k_d(r)\,dr/r$.
\begin{enumerate}
\item \emph{(The exponent grows at most like $c\log|\xi|$.)} For every $\xi$ with $|\xi| \ge 1$,
\[
  F_d(\xi) \le c\log|\xi| + \tfrac{13}{4}\,c + 2C_k ,
  \qquad\text{so}\qquad
  e^{-F_d(\xi)} \ge e^{-13c/4 - 2C_k}\,|\xi|^{-c} .
\]
\item \emph{(The exponent is logarithmic.)} For every $a > 0$,
$F_d(a\xi) - F_d(\xi) \to c\log a$ as $|\xi| \to \infty$. Equivalently, for a unit vector $e$ the
function $g(u) := F_d(\sqrt u\,e)$, $u \ge 0$, satisfies $g(au) - g(u) \to \tfrac c2\log a$ as
$u \to \infty$, for every $a > 0$.
\end{enumerate}
\end{lemma}

Clause (1) is the lower bound on the transform quoted in \S\ref{sec:isotropic-origin}, and it
is used in Theorem~\ref{thm:origin-boundedness-d}(4) and in
Corollary~\ref{cor:origin-profile-exponent}(2). For a finite $c$, clause (2) says more than that
corollary: $e^{-F_d(\xi)} = |\xi|^{-c}L(|\xi|)$ with $L$ slowly varying, and
Proposition~\ref{prop:isotropic-exponent-profile} gives $L$ through the polar profile.

Two bounds complement the rate of \S\ref{sec:isotropic-origin}. Each is for one kernel, with
its constants. The first bounds the spike from below, for $d \ge 2$ and a finite $c$, and is
stated for the kernel as a measure before it is stated for its density.

% shared with blueprint prop:origin-lower-bound-d+69c7d8a91a64
\begin{proposition}[a pointwise lower bound at the origin]\label{prop:origin-lower-bound-d}
Let $d \ge 2$, let $(0, k_d)$ be an isotropic origin datum in dimension $d$ with Gaussian
coefficient $0$, let $\mu$ be its law, and suppose that $c := k_d(0{+})$ is finite, $0 < c < \infty$.
Then there are $\delta > 0$ and $c_1 > 0$ with
\begin{equation}\label{eq:origin-lower-bound}
  \mu(E) \ge c_1\int_E |x|^{c-d}\,dx \qquad\text{for every Borel set } E \subseteq \{|x| < \delta\} .
\end{equation}
In particular:
\begin{enumerate}
\item the density $q$ of the absolutely continuous part of $\mu$ satisfies
$q(x) \ge c_1|x|^{c-d}$ for Lebesgue-almost every $x$ with $|x| < \delta$; so does every density of
$\mu$ when $\mu$ has one, and every representative of $q$ is at least
$c_1|x|^{c-d}$ almost everywhere near the origin, which tends to $\infty$ as $x \to 0$ when $c < d$;
\item $\mu(\{|x| < r\}) \ge c_1|S^{d-1}|\,c^{-1}\,r^{c}$ for $0 < r \le \delta$.
\end{enumerate}
The constants can be taken as follows. Let $\varepsilon_0 \in (0,1]$ satisfy
$k_d(\varepsilon_0) \ge c/2$, put $C_k := \int_1^\infty k_d(r)\,dr/r$, let $V_d$ be the volume of the
unit ball, and put
\[
  \tau := \min\Bigl\{\tfrac12,\ \bigl(\tfrac{3}{4d^2c}\bigr)^{1/2}\Bigr\},\quad
  \kappa_0 := \Bigl(\frac{c}{2|S^{d-1}|}\Bigr)^{2}\frac{V_d}{8^d},\quad
  \delta := \frac{\varepsilon_0}{2},\quad
  c_1 := \frac{\kappa_0}{4}\,\tau^{c}e^{-C_k}\min\{2^{c-d},4^{c-d}\} .
\]
\end{proposition}

For $c \le d$ the power cannot be improved. A bound $c_1|x|^{\beta-d}$ with $\beta < c$ would give
$\int g_t\,d\mu \ge C\,t^{-(d-\beta)/2}$, against the upper bound of
Theorem~\ref{thm:origin-boundedness-d}(4) at a power between $\beta$ and $c$. For $c > d$ the
bound tends to $0$ and says little, the density being continuous and positive at the centre.
The proof uses $d \ge 2$ at one step, the choice of a nonzero vector orthogonal to a given
point. On the line, where that step fails, the rate gives the bound near the centre for
$c < 1$, without a radius.

The second proposition bounds the spike from above, when $c \le d$, at every power below $c$.
The upper bound of clause (4) of Theorem~\ref{thm:origin-boundedness-d} on the blurred kernel
gives a bound on the mass of a ball, and a radially nonincreasing density is at most its
average over a ball.

% shared with blueprint prop:origin-upper-bound-d+0f6cbfb1cb1f
\begin{proposition}[an upper bound at the origin, and the order of the density there]\label{prop:origin-upper-bound-d}
Let $d \ge 1$, let $(0, k_d)$ be an isotropic origin datum in dimension $d$ with Gaussian
coefficient $0$, let $\mu$ be its law, and suppose that $c := k_d(0{+})$ satisfies $0 < c \le d$.
Let $V_d$ be the volume of the unit ball of $\RR^d$.
\begin{enumerate}
\item \emph{(The mass of a ball.)} For every $\beta$ with $0 < \beta < c$ there is
$A_\beta < \infty$ with
\[
  \mu(\{|x| < r\}) \le A_\beta\,r^{\beta} \qquad\text{for every } r > 0 .
\]
One may take $A_\beta = (2\pi)^{d/2}e^{1/2}\,C'_\beta$, for every constant $C'_\beta$ with which
the upper bound of Theorem~\ref{thm:origin-boundedness-d}(4) holds.
\item \emph{(The density.)} The law $\mu$ has a density $p$ with $p(x) = q(|x|)$ for $x \ne 0$,
where $q : (0,\infty) \to [0,\infty)$ is nonincreasing. For every such $q$, every $0 < \beta < c$, and
$A_\beta = (2\pi)^{d/2}e^{1/2}\,C'_\beta$ with $C'_\beta$ as in clause (1),
\[
  q(r) \le \frac{A_\beta}{V_d}\,r^{\beta-d} \qquad\text{for every } r > 0 .
\]
\item \emph{(The order at the origin, for $d \ge 2$.)} Let $d \ge 2$, let $\delta$ and $c_1$ be
the constants displayed at the end of Proposition~\ref{prop:origin-lower-bound-d}, and let $q$ be as in
clause (2). Then
$q(r) \ge c_1\,r^{c-d}$ for every $r$ with $0 < r < \delta$, with no exceptional set. With
$p(x) = q(|x|)$,
\[
  \frac{\log p(x)}{\log(1/|x|)} \to d - c \quad (x \to 0), \qquad
  \frac{\log \mu(\{|x| < r\})}{\log r} \to c \quad (r \downarrow 0) .
\]
\end{enumerate}
\end{proposition}

With the lower bound, clause (3) gives the logarithmic order of the density for $d \ge 2$. At
$c = d$ its first limit is $0$, and how the density grows there is clause (3) of
Proposition~\ref{prop:origin-rate-profile-d}. \suppl{} gives the constant $A_\beta$ in closed
form, after the proof.

Proposition~\ref{prop:profile-one-dimension-down} rests on the next lemma: the part of a
rotation-invariant \Levy{} measure on $\RR^{d+1}$ that a projection sees in $\RR^d$, its
\emph{shadow}, has a polar profile in closed form, and that profile is continuous.

% shared with blueprint lem:shadow-polar-profile+1355442004bb
\begin{lemma}[the polar profile of a shadow]\label{lem:shadow-polar-profile}
Let $d \ge 1$, let $(a, k')$ be an isotropic origin datum in dimension $d + 1$, with \Levy{}
measure $\nu'$ and exponent $F'$, and let $\Pi(x,y) := x$ be the projection of
$\RR^{d+1} = \RR^d \times \RR$ on the first $d$ coordinates. Put
\[
  C_d := \frac{|S^{d-1}|}{|S^{d}|}, \qquad
  \bar k(r) := C_d\int_{\RR} k'\bigl(r\sqrt{1+u^2}\bigr)\,(1+u^2)^{-(d+1)/2}\,du \quad (r > 0) .
\]
\begin{enumerate}
\item $C_d\int_{\RR}(1+u^2)^{-(d+1)/2}\,du = 1$. So $\bar k(r)$ is a mean of the values of $k'$
on $[r,\infty)$, and $0 \le \bar k(r) \le k'(r)$.
\item $\bar k$ is nonincreasing and continuous on $(0,\infty)$, and $\bar k(0{+}) = k'(0{+})$ in
$[0,\infty]$.
\item $(a, \bar k)$ is an isotropic origin datum in dimension $d$. Its \Levy{} measure $\bar\nu$
is the image of $\nu'$ under $\Pi$, restricted to $\RR^d \setminus \{0\}$: for every nonnegative
Borel $g$ on $\RR^d$ vanishing at the origin, $\int g(\Pi z)\,\nu'(dz) = \int g\,d\bar\nu$. The
restriction removes nothing: the image of $\nu'$ gives the origin of $\RR^d$ the mass
$\nu'(\{z \ne 0 : \Pi z = 0\}) = 0$. The exponent of $(a,\bar k)$ is
$\xi \mapsto F'\bigl((\xi,0)\bigr)$, $\xi \in \RR^d$.
\end{enumerate}
In particular, if $F'(\zeta) = F(|\zeta|)$ for a function $F$ on $[0,\infty)$, then $(a,\bar k)$
is an isotropic origin datum in dimension $d$ with exponent $\xi \mapsto F(|\xi|)$.
\end{lemma}

Applied $n$ times, the lemma gives the shadow in dimension $d$ of a datum in dimension $d + n$.
The last lemma says that the bridge formula and the polar profile \ref{eq:polar-profile-d} hold
for a delay datum whose delay profile is not monotone, as that of the separating family is not.

% Q-0406 (the author's review of section 3, item 5): the paper's copy says "a delay datum in the
% sense of section 2.3, with exponent F_I" and drops the conditions and the formula that the
% blueprint node restates from [B, Def. 4.11] after "restated:", which section 2.3 gives; it keeps
% the one point the lemma turns on, that the delay profile need not be monotone. Deliberately
% edited, and pinned.
% shared with blueprint lem:delay-polar-profile@fa344f467980
\begin{lemma}[the polar profile of a delay datum, without monotonicity]\label{lem:delay-polar-profile}
Let $d \ge 1$ and let $(b_0, k_{\mathrm I})$ be a delay datum in the sense of \restated{B}{Def.~4.11}{sec:line-bridge},
with exponent $F_{\mathrm I}$; its delay profile $k_{\mathrm I}$ is measurable and not necessarily
monotone. Let $k_d$ be its polar profile \ref{eq:polar-profile-d}. Then $k_d$ is finite and measurable on
$(0,\infty)$, with
\[
  \int_0^1 k_d(r)\,r\,dr < \infty
  \qquad\text{and}\qquad
  \int_1^\infty k_d(r)\,\frac{dr}{r} < \infty ,
\]
and
\[
  F_{\mathrm I}\bigl(\tfrac12|\xi|^2\bigr) = \tfrac12 b_0|\xi|^2
   + \int_{\RR^d}\bigl(1 - \cos\langle\xi,x\rangle\bigr)\,\frac{k_d(|x|)}{|S^{d-1}|\,|x|^d}\,dx ,
   \qquad \xi \in \RR^d .
\]
In particular, if $k_d$ is nonincreasing, then $(b_0,k_d)$ is an isotropic origin datum in
dimension $d$ whose exponent is $\xi \mapsto F_{\mathrm I}(\tfrac12|\xi|^2)$.
\end{lemma}

The last clause of the lemma is also what makes the datum of the \Matern{} kernel in
Example~\ref{ex:matern} a subordinated one.

\subsection{The statements proved in print}\label{app:print-proofs}

% Q-0448: the statements of the second list of section 1.1, each with its proof. A statement
% that Lean comes to check leaves this subsection for supp-proofs.tex, and the list.
This subsection has the proofs of the statements that are not machine-checked
(\S\ref{sec:trust-base}). Three of them are stated in the main text,
Lemma~\ref{lem:exponent-profile-constant} and Propositions~\ref{prop:radial-dictionary}
and~\ref{prop:step-exponent-dimensions}, and two are stated here,
Corollary~\ref{cor:dictionary-filtration} and
Proposition~\ref{prop:separating-family-mixtures}. The proof of
Proposition~\ref{prop:profile-one-dimension-down}, which Lean checked, is given with them,
because the corollary reads that proposition in the dictionary.

\begin{proof}[Proof of Lemma~\ref{lem:exponent-profile-constant}]
\emph{Three facts about $\Gamma$.} For $a > 0$, $\Gamma(a) = \int_0^\infty t^{a-1}e^{-t}\,dt$. On
a compact interval $[a_0,a_1] \subset (0,\infty)$ the derivative of the integrand in $a$,
$t^{a-1}e^{-t}\log t$, is bounded in absolute value by
$(t^{a_0-1} + t^{a_1-1})\,e^{-t}\,|\log t|$, which is integrable on $(0,\infty)$. So one may
differentiate under the integral sign:
\[
  \Gamma'(a) = \int_0^\infty t^{a-1}e^{-t}\log t\,dt ,
\]
the integral converging absolutely. Consequently, if a random variable $T$ has the gamma law of
shape $a$, with density $t^{a-1}e^{-t}/\Gamma(a)$ on $(0,\infty)$, then $\log T$ is integrable and
$\mathbb{E}\log T = \psi(a)$. At $a = 1$ this reads
$\gamma_{\mathrm E} = -\int_0^\infty e^{-v}\log v\,dv$. Next, $\Gamma(a+1) = a\,\Gamma(a)$ gives
$\psi(a+1) = \psi(a) + 1/a$. Last, the duplication formula
$\Gamma(a)\,\Gamma(a+\tfrac12) = 2^{1-2a}\sqrt\pi\,\Gamma(2a)$ gives, on taking logarithms and
differentiating, $\psi(a) + \psi(a+\tfrac12) = 2\psi(2a) - 2\log 2$. At $a = \tfrac12$,
\[
  \psi(\tfrac12) = \psi(1) - 2\log 2 = -\gamma_{\mathrm E} - 2\log 2 .
\]

\emph{Clause (1).} Let $G(X) = \Cin(X) - \log^+ X$ be as in the proof of
Proposition~\ref{prop:isotropic-exponent-profile}, so that $-3 \le G \le \tfrac{13}{4}$ and
$G(X) \to g_\infty$ as $X \to \infty$. For $0 < \varepsilon < 1$ put
\[
  J(\varepsilon) := \int_0^\infty \bigl(1 - \cos u\bigr)\,e^{-\varepsilon u}\,\frac{du}{u} .
\]
\emph{(a)} $(1 - \cos u)/u = \int_0^1 \sin(tu)\,dt$, and
$\int_0^\infty \sin(tu)\,e^{-\varepsilon u}\,du = t/(\varepsilon^2 + t^2)$. The function
$|\sin(tu)|\,e^{-\varepsilon u}$ is integrable on $(0,1) \times (0,\infty)$, so Fubini's theorem
gives
\[
  J(\varepsilon) = \int_0^1 \frac{t\,dt}{\varepsilon^2 + t^2}
   = \tfrac12\log\bigl(1 + \varepsilon^{-2}\bigr)
   = \log\frac1\varepsilon + \tfrac12\log\bigl(1 + \varepsilon^{2}\bigr) .
\]
\emph{(b)} $e^{-\varepsilon u} = \int_u^\infty \varepsilon e^{-\varepsilon x}\,dx$. The integrand
$(1 - \cos u)\,u^{-1}\,\varepsilon e^{-\varepsilon x}\,\mathbf 1\{u < x\}$ is nonnegative, so
Tonelli's theorem gives
\[
  J(\varepsilon) = \int_0^\infty \varepsilon e^{-\varepsilon x}\,\Cin(x)\,dx
   = \int_0^\infty e^{-v}\,\Cin(v/\varepsilon)\,dv .
\]
\emph{(c)} For $v > 0$,
\[
  \Cin(v/\varepsilon) - \log\frac1\varepsilon
   = G(v/\varepsilon) + \Bigl[\log^+\frac v\varepsilon - \log\frac1\varepsilon\Bigr] .
\]
The bracket equals $\log v$ when $v \ge \varepsilon$. When $v < \varepsilon$ it equals
$\log\varepsilon$, which lies between $\log v$ and $0$. So the bracket is bounded by $|\log v|$ in
absolute value and tends to $\log v$ as $\varepsilon \downarrow 0$. As $|G| \le \tfrac{13}{4}$ and
$G(v/\varepsilon) \to g_\infty$, dominated convergence with the integrable function
$e^{-v}(\tfrac{13}{4} + |\log v|)$ gives, by (b) and $\int_0^\infty e^{-v}dv = 1$,
\[
  J(\varepsilon) - \log\frac1\varepsilon \to g_\infty + \int_0^\infty e^{-v}\log v\,dv
   = g_\infty - \gamma_{\mathrm E} \qquad (\varepsilon \downarrow 0) .
\]
\emph{(d)} By (a) the left side is $\tfrac12\log(1 + \varepsilon^2)$, which tends to $0$. Hence
$g_\infty = \gamma_{\mathrm E}$.

\emph{Clause (2).} Let $Z$ be a standard Gaussian vector of $\RR^d$, with density
$(2\pi)^{-d/2}e^{-|x|^2/2}$, and $e$ a unit vector. In polar coordinates, for a Borel
$h \ge 0$ on $S^{d-1}$ and a Borel $f \ge 0$ on $(0,\infty)$,
\[
  \mathbb{E}\bigl[h(Z/|Z|)\,f(|Z|)\bigr]
   = \int_{S^{d-1}}h\,d\sigma_{d-1}\int_0^\infty f(r)\,(2\pi)^{-d/2}e^{-r^2/2}\,r^{d-1}\,dr .
\]
With $h = 1$ and the substitution $t = r^2/2$, using $|S^{d-1}| = 2\pi^{d/2}/\Gamma(d/2)$, this
shows that $|Z|^2/2$ has the gamma law of shape $d/2$. With $f = 1$ it shows that $Z/|Z|$ has the
law $\sigma_{d-1}/|S^{d-1}|$. The coordinate $\langle e,Z\rangle$ is a standard Gaussian variable,
so $\langle e,Z\rangle^2/2$ has the gamma law of shape $\tfrac12$, by the case $d = 1$. By the
first fact, $\log|Z|$ and $\log|\langle e,Z\rangle|$ are integrable, with
\[
  \mathbb{E}\log|Z| = \tfrac12\bigl[\log 2 + \psi(\tfrac d2)\bigr], \qquad
  \mathbb{E}\log|\langle e,Z\rangle| = \tfrac12\bigl[\log 2 + \psi(\tfrac12)\bigr] .
\]
Since $\langle e,Z\rangle = |Z|\,\langle e, Z/|Z|\rangle$,
\[
  \lambda_d = \mathbb{E}\log\bigl|\langle e, Z/|Z|\rangle\bigr|
   = \mathbb{E}\log|\langle e,Z\rangle| - \mathbb{E}\log|Z|
   = \tfrac12\bigl[\psi(\tfrac12) - \psi(\tfrac d2)\bigr] .
\]

\emph{Clause (3).} The first expression is clauses (1) and (2). The second follows from
$\psi(\tfrac12) = -\gamma_{\mathrm E} - 2\log 2$. At $d = 1$ the first expression is
$\gamma_{\mathrm E}$. At $d = 2$ the second is
$\tfrac12(\gamma_{\mathrm E} - \psi(1)) - \log 2 = \gamma_{\mathrm E} - \log 2$. By
$\psi(a+1) = \psi(a) + 1/a$ at $a = d/2$,
$\varkappa_{d+2} = \varkappa_d - \tfrac12\cdot\tfrac2d = \varkappa_d - 1/d$, and so
$\varkappa_3 = \varkappa_1 - 1$.

\emph{Clause (4).} By the second expression of clause (3),
$-c\varkappa_d = c\log 2 - \tfrac c2(\gamma_{\mathrm E} - \psi(\tfrac d2))$, which is the first
display. The second follows from
$R(d,c) = \Gamma(\tfrac{d-c}{2})/(2^{c}\pi^{d/2}\Gamma(\tfrac c2))$. For the third take $c = d$
in the first display and use $|S^{d-1}| = 2\pi^{d/2}/\Gamma(\tfrac d2)$:
$(2\pi)^{-d}\cdot 2\pi^{d/2}\,\Gamma(\tfrac d2)^{-1}\cdot 2^{d} = 2/(\pi^{d/2}\Gamma(\tfrac d2))$.
\end{proof}

\begin{proof}[Proof of Proposition~\ref{prop:radial-dictionary}]
\emph{The function $\Omega_d$.} For unit vectors $e, e'$ take $Q \in \Orth(d)$ with $Qe = e'$.
Then $\langle e',\eta\rangle = \langle e,Q^{\tp}\eta\rangle$, and $\sigma_{d-1}$ is invariant
under $Q^{\tp}$, so the integral defining $\Omega_d(t)$ is the same for $e'$ as for $e$. At
$d = 1$ the sphere is $\{\pm1\}$ and $\Omega_1(t) = \cos t$. $\Omega_d$ is continuous, by
dominated convergence, $|\Omega_d| \le 1$ and $\Omega_d(0) = 1$. For $x \in \RR^d$ and $t \ge 0$,
\begin{equation}\label{eq:sphere-mean-cos}
  \frac{1}{|S^{d-1}|}\int_{S^{d-1}}e^{it\langle x,\theta\rangle}\,\sigma_{d-1}(d\theta)
   = \Omega_d\bigl(t|x|\bigr) .
\end{equation}
Indeed the imaginary part vanishes by the symmetry $\theta \mapsto -\theta$ of $\sigma_{d-1}$.
The real part is the definition of $\Omega_d(t|x|)$ with $e = x/|x|$ when $x \ne 0$, and both
sides are $1$ when $x = 0$.

\emph{The preamble.} Let $X$ be rotation invariant. As $-I \in \Orth(d)$, $-X$ has the law of
$X$, so the characteristic function is real:
$\mathbb{E}\,e^{i\langle\xi,X\rangle} = \mathbb{E}\cos\langle\xi,X\rangle$. Let $\xi \in \RR^d$
and let $\theta$ be a unit vector. There is $Q \in \Orth(d)$ with $Q\xi = |\xi|\theta$, and
$Q^{\tp}X$ has the law of $X$, so
\[
  \mathbb{E}\cos\langle\xi,X\rangle = \mathbb{E}\cos\langle\xi,Q^{\tp}X\rangle
   = \mathbb{E}\cos\bigl(|\xi|\langle\theta,X\rangle\bigr) .
\]
The left side does not depend on $\theta$. Averaging over $\theta$ against
$\sigma_{d-1}/|S^{d-1}|$, exchanging the two integrals of a bounded continuous function, and using
\ref{eq:sphere-mean-cos} gives
$\mathbb{E}\,e^{i\langle\xi,X\rangle} = \mathbb{E}\,\Omega_d(|\xi||X|) = \phi(|\xi|)$. By the
uniqueness theorem for characteristic functions, $\phi$ determines the law of $X$. Let $R$ have
the law of $|X|$ and let $U$ be uniform on $S^{d-1}$, that is of law $\sigma_{d-1}/|S^{d-1}|$,
and independent of $R$. By \ref{eq:sphere-mean-cos} and Fubini's theorem,
$\mathbb{E}\,e^{i\langle\xi,RU\rangle} = \mathbb{E}\,\Omega_d(R|\xi|) = \phi(|\xi|)$. So $RU$ has
the law of $X$, by the same theorem.

\emph{Step A: the identity.} Let $R$ be a random variable with values in $(0,\infty)$, let
$0 < c < \infty$, and put $k(r) := c\,\Pr(R > r)$ for $r > 0$ and
$\phi(v) := \mathbb{E}\,\Omega_d(vR)$ for $v \ge 0$. Then for every $\rho > 0$,
\begin{equation}\label{eq:dictionary-identity}
  \int_0^\infty\bigl(1 - \Omega_d(\rho r)\bigr)\,k(r)\,\frac{dr}{r}
   = c\int_0^\rho\bigl(1 - \phi(v)\bigr)\,\frac{dv}{v} ,
\end{equation}
both sides in $[0,\infty]$. Indeed, as $1 - \Omega_d \ge 0$, Tonelli's theorem applies twice, with
the substitution $r = vR/\rho$ in between:
\[
\begin{aligned}
  \int_0^\infty\bigl(1 - \Omega_d(\rho r)\bigr)\,c\,\Pr(R > r)\,\frac{dr}{r}
   &= c\,\mathbb{E}\int_0^R\bigl(1 - \Omega_d(\rho r)\bigr)\frac{dr}{r}
   = c\,\mathbb{E}\int_0^\rho\bigl(1 - \Omega_d(vR)\bigr)\frac{dv}{v} \\
   &= c\int_0^\rho\bigl(1 - \phi(v)\bigr)\frac{dv}{v} .
\end{aligned}
\]
Moreover, if $\nu$ is the measure on $\RR^d \setminus \{0\}$ with the density
$k(|x|)\,(|S^{d-1}|\,|x|^d)^{-1}$, then for $\xi \ne 0$, by polar coordinates, Tonelli's theorem
and \ref{eq:sphere-mean-cos},
\begin{equation}\label{eq:dictionary-polar}
  \int_{\RR^d}\bigl(1 - \cos\langle\xi,x\rangle\bigr)\,\nu(dx)
   = \int_0^\infty\bigl(1 - \Omega_d(|\xi| r)\bigr)\,k(r)\,\frac{dr}{r} .
\end{equation}
The same holds for every nonnegative Borel $k$ in place of this one.

\emph{Step B: the two integrability conditions.} With $R$, $c$ and $k$ as in Step~A, Tonelli's
theorem gives
\[
  \int_1^\infty k(r)\,\frac{dr}{r} = c\,\mathbb{E}\int_1^{\max\{R,1\}}\frac{dr}{r}
   = c\,\mathbb{E}\log^+R ,
\]
and $\int_0^1 k(r)\,r\,dr \le c/2$ because $k \le c$.

\emph{Step C: the logarithmic moment read on $\phi$.} Put
$H_d(s) := \int_0^s (1 - \Omega_d(v))\,dv/v$ for $s \ge 0$. It is nonnegative and nondecreasing.
The pair $(0,\mathbf 1_{(0,1)})$ is an isotropic origin datum in dimension $d$, with profile at
the origin $1$ and with $\int_1^\infty\mathbf 1_{(0,1)}(r)\,dr/r = 0$. By
\ref{eq:dictionary-polar} and \ref{eq:dictionary-identity} with $R \equiv 1$ and $c = 1$, its
exponent at a point of norm $s$ is $H_d(s)$. So Lemma~\ref{lem:isotropic-exponent-log}(1) gives
$H_d(s) \le \log s + \tfrac{13}{4}$ for $s \ge 1$, and hence
$H_d(s) \le \log^+ s + \tfrac{13}{4}$ for every $s \ge 0$, $H_d$ being nondecreasing. And
Corollary~\ref{cor:origin-profile-exponent}(2) gives $H_d(s)/\log s \to 1$ as $s \to \infty$, so there is
$s_0 > 1$ with $H_d(s) \ge \tfrac12\log s$ for $s \ge s_0$. Now let $R$ and $\phi$ be as in
Step~A. By Tonelli's theorem and the substitution $v \mapsto vR$,
\[
  \int_0^1\bigl(1 - \phi(v)\bigr)\frac{dv}{v}
   = \mathbb{E}\int_0^1\bigl(1 - \Omega_d(vR)\bigr)\frac{dv}{v} = \mathbb{E}\,H_d(R) .
\]
By the two bounds on $H_d$,
$\tfrac12\,\mathbb{E}\log^+R - \tfrac12\log s_0 \le \mathbb{E}\,H_d(R) \le
\mathbb{E}\log^+R + \tfrac{13}{4}$. So $\int_0^1(1 - \phi(v))\,dv/v$ is finite if and only if
$\mathbb{E}\log^+R$ is.

\emph{Clause (1).} Let $X$ be as in the clause and $R := |X|$, which is positive almost surely.
By the preamble $\phi(v) = \mathbb{E}\,\Omega_d(vR)$, and the equivalence is Step~C. Suppose
$\mathbb{E}\log^+|X| < \infty$ and put $k_d(r) := c\,\Pr(|X| > r)$. It is nonincreasing with
values in $[0,c]$, and $k_d(0{+}) = c\,\Pr(|X| > 0) = c$. By Step~B it satisfies the two
conditions of Definition~\ref{def:isotropic-origin-datum}. So $(0,k_d)$ is an isotropic origin datum in
dimension $d$. By \ref{eq:dictionary-polar} and \ref{eq:dictionary-identity} its exponent at
$\xi \ne 0$ is $c\int_0^{|\xi|}(1 - \phi(v))\,dv/v = F(|\xi|)$, and both vanish at $\xi = 0$. By
Proposition~\ref{prop:isotropic-polar-reading}(2), $F \in \mathcal{A}_d$. By
Corollary~\ref{cor:origin-profile-exponent}(2) for this datum, read along a ray, $F(r)/\log r \to c$.

\emph{Clause (2), existence.} Let $F \in \mathcal{A}_d$ with $F(r)/\log r \to c \in (0,\infty)$.
By Proposition~\ref{prop:isotropic-polar-reading}(2), $\xi \mapsto F(|\xi|)$ is the exponent of an
isotropic origin datum $(a,k_d)$ in dimension $d$. By Corollary~\ref{cor:origin-profile-exponent}(1),
$F(r)/r^2 \to \tfrac12 a$; as $F(r)/\log r$ has a finite limit, $F(r)/r^2 \to 0$, so $a = 0$.
Then Corollary~\ref{cor:origin-profile-exponent}(2) gives $F(r)/\log r \to k_d(0{+})$, so $k_d(0{+}) = c$.
As in the proof of Lemma~\ref{lem:isotropic-exponent-log}, $k_d(r) \to 0$ as $r \to \infty$. Put
$G(r) := k_d(r{+})/c$ for $r > 0$. It is nonincreasing and right-continuous, it tends to $1$ as
$r \downarrow 0$ and to $0$ as $r \to \infty$, and $c\,G = k_d$ off the countable set of the
discontinuities of $k_d$. So there is a random variable $R$ with values in $(0,\infty)$ and
$\Pr(R > r) = G(r)$ for every $r > 0$. Let $U$ be uniform on $S^{d-1}$ and independent of $R$,
and $X := RU$. Then $X$ is rotation invariant, $|X| = R$, $\Pr(X = 0) = 0$, and by Step~B and
$c\,G = k_d$ almost everywhere,
$c\,\mathbb{E}\log^+|X| = \int_1^\infty k_d(r{+})\,dr/r = \int_1^\infty k_d(r)\,dr/r < \infty$.
Its function $\phi$ is that of
Step~A. Let $\rho > 0$ and let $e$ be a unit vector. By \ref{eq:dictionary-polar} for $k_d$, the
equality $k_d = c\,\Pr(R > \cdot)$ almost everywhere, and \ref{eq:dictionary-identity},
\[
  F(\rho) = \int_0^\infty\bigl(1 - \Omega_d(\rho r)\bigr)\,k_d(r)\,\frac{dr}{r}
   = c\int_0^\rho\bigl(1 - \phi(v)\bigr)\frac{dv}{v} .
\]
So $F$ is the function of clause (1) for $c$ and $X$. As $\phi$ is continuous, the integrand
$(1 - \phi(v))/v$ is continuous on $(0,\infty)$, so $F$ is continuously differentiable there, with
$rF'(r) = c\,(1 - \phi(r))$. The datum $(0,k_d)$ was any datum with exponent $F(|\cdot|)$, and
$\Pr(|X| > r) = k_d(r{+})/c$ by construction.

\emph{Clause (2), uniqueness.} Suppose $F$ is also the function of clause (1) for $c_1$ and
$X_1$, with $\mathbb{E}\log^+|X_1| < \infty$ and radial function $\phi_1$. By clause (1),
$F(r)/\log r \to c_1$, so $c_1 = c$. Both integrands being continuous, differentiation gives
$c\,(1 - \phi_1(r)) = rF'(r) = c\,(1 - \phi(r))$ for every $r > 0$, and
$\phi_1(0) = \phi(0) = 1$. So $\phi_1 = \phi$, and $X_1$ has the law of $X$ by the preamble. In
particular $\Pr(|X| > r) = k_d(r{+})/c$ holds for the one law and every datum, since the
construction above can start from any datum. \emph{The bijection.} By clause (1) the map sends a
pair $(c, \text{law of } X)$ into the set named; by the existence part it is onto, and by the
uniqueness part it is one to one. The set is that of the members of $\mathcal{A}_d$ with
Gaussian coefficient $0$ and profile at the origin in $(0,\infty)$: one inclusion was shown at
the start of the existence part, and the other is Corollary~\ref{cor:origin-profile-exponent}(2).

\emph{Clause (3), if.} Let $X'$ be a random vector in $\RR^{d+1}$ with characteristic function
$\zeta \mapsto \phi(|\zeta|)$. For $Q \in \Orth(d+1)$ the characteristic function of $QX'$ is
$\zeta \mapsto \phi(|Q^{\tp}\zeta|) = \phi(|\zeta|)$, so $QX'$ has the law of $X'$: $X'$ is
rotation invariant. Let $\Pi$ be the projection on the first $d$ coordinates. The characteristic
function of $\Pi X'$ at $\xi \in \RR^d$ is that of $X'$ at $(\xi,0)$, that is $\phi(|\xi|)$. So
$\Pi X'$ has the law of $X$, and $\Pr(X' = 0) \le \Pr(\Pi X' = 0) = \Pr(X = 0) = 0$. By the
preamble in dimension $d + 1$, the radial function of $X'$ is $\phi$. As
$\int_0^1(1 - \phi(u))\,du/u = F(1)/c < \infty$, clause (1) in dimension $d + 1$ gives
$\mathbb{E}\log^+|X'| < \infty$ and $F \in \mathcal{A}_{d+1}$, with $F$ the function of that
clause for $c$ and $X'$.

\emph{Clause (3), only if.} Let $F \in \mathcal{A}_{d+1}$. By clause (1) in dimension $d$,
$F(r)/\log r \to c$. So clause (2) applies in dimension $d + 1$: there is a rotation-invariant
$X'$ in $\RR^{d+1}$ with $\Pr(X' = 0) = 0$ and $\mathbb{E}\log^+|X'| < \infty$ whose radial
function is $r \mapsto 1 - rF'(r)/c$ on $(0,\infty)$. By clause (2) in dimension $d$ this
function is $\phi$. Both radial functions equal $1$ at $0$. So the characteristic function of
$X'$ is $\zeta \mapsto \phi(|\zeta|)$.
\end{proof}

\begin{proof}[Proof of Proposition~\ref{prop:profile-one-dimension-down}]
\emph{Clause (1).} By Lemma~\ref{lem:shadow-polar-profile}, with $\bar k$ the function displayed in
the clause, $(a',\bar k)$ is an isotropic origin datum in dimension $d$ with exponent
$\xi \mapsto F(|\xi|)$. So $F \in \mathcal{A}_d$, by
Proposition~\ref{prop:isotropic-polar-reading}(2). Let $(a,k_d)$ be any isotropic origin datum in
dimension $d$ with that exponent. By the uniqueness in Proposition~\ref{prop:isotropic-polar-reading}(1),
$a = a'$ and $k_d = \bar k$ almost everywhere on $(0,\infty)$. Let $r > 0$, and choose points
$r_n \downarrow r$ and $s_n \uparrow r$ at which $k_d = \bar k$. As $k_d$ is nonincreasing,
$\bar k(r_n) = k_d(r_n) \le k_d(r) \le k_d(s_n) = \bar k(s_n)$, and both ends tend to
$\bar k(r)$, because $\bar k$ is continuous (Lemma~\ref{lem:shadow-polar-profile}(2)). So
$k_d(r) = \bar k(r)$ for every $r > 0$, and $k_d$ is continuous. The next sentence of the clause
is the contrapositive.

\emph{Clause (2).} The function $k := b\,\mathbf 1_{(0,\tau)}$ is nonnegative and nonincreasing
on $(0,\infty)$, with $\int_0^1 k(r)\,r\,dr \le b/2$ and
$\int_1^\infty k(r)\,dr/r = b\log^+\tau$. So $(0,k)$ is an isotropic origin datum in dimension
$d$ (Definition~\ref{def:isotropic-origin-datum}). Its exponent $F_d$ is finite and nonnegative, as that
of every datum, and radial: for
$Q \in \Orth(d)$ the substitution $x = Qy$ in \ref{eq:isotropic-exponent} gives
$F_d(Q\xi) = F_d(\xi)$, because $Q$ preserves Lebesgue measure and the radial density of the
\Levy{} measure. So $F_d(\xi) = G(|\xi|)$ with $G(r) := F_d(re_1)$ for $r \ge 0$, and
$G \in \mathcal{A}_d$ by Proposition~\ref{prop:isotropic-polar-reading}(2). The profile $k$ has the
left limit $b$ and the right limit $0$ at $\tau$. By clause (1), $G \notin \mathcal{A}_{d+1}$.
So $\mathcal{A}_{d+1} \ne \mathcal{A}_d$. The inclusion
$\mathcal{A}_{d+1} \subseteq \mathcal{A}_d$ follows from clause (1), since for every
$F \in \mathcal{A}_{d+1}$, $\zeta \mapsto F(|\zeta|)$ is the exponent of an isotropic origin datum
in dimension $d + 1$, by Proposition~\ref{prop:isotropic-polar-reading}(2).

\emph{Clause (3).} At $d = 1$ the sphere is $S^0 = \{\pm1\}$, with $\sigma_0$ the counting
measure. By \ref{eq:polar-decomposition} the exponent of the datum
$(0, b\,\mathbf 1_{(0,\tau)})$ at $\omega \in \RR$ is
$b\int_0^\tau(1 - \cos\omega r)\,dr/r$, which is $b\,\Cin(\tau|\omega|)$: for $\omega \ne 0$ by
the substitution $v = |\omega| r$, and both vanish at $\omega = 0$. So
$G(r) = b\,\Cin(\tau r)$.
This is the $\Cin$ ray of Proposition~\ref{prop:line-extreme-rays}, the admissible exponent with
Gaussian coefficient $0$ and profile $b\,\mathbf 1_{(0,\tau)}$. By clause (2) it is not in
$\mathcal{A}_2$. By clause (1), with the datum that
Proposition~\ref{prop:isotropic-polar-reading}(2) supplies,
$\mathcal{A}_{D+1} \subseteq \mathcal{A}_D$ for every $D \ge 1$, so
$\mathcal{A}_D \subseteq \mathcal{A}_2$ for $D \ge 2$, and the ray is in none of them.
\end{proof}

The corollary reads the first two clauses of
Proposition~\ref{prop:profile-one-dimension-down} in the dictionary of
Proposition~\ref{prop:radial-dictionary}.

% shared with blueprint cor:dictionary-filtration+60d4f00473ee
\begin{corollary}[the test and the step datum, read in the dictionary]\label{cor:dictionary-filtration}
Let $d \ge 1$, and let $\Omega_d$ be the function of Proposition~\ref{prop:radial-dictionary}.
\begin{enumerate}
\item \emph{(The jump size has no atom.)} In the terms of
Proposition~\ref{prop:radial-dictionary}: if $F$ and $X$ are as in its clause (1), with
$\mathbb{E}\log^+|X| < \infty$, and $F \in \mathcal{A}_{d+1}$, then $\Pr(|X| = r) = 0$ for every
$r > 0$.
\item \emph{(The step datum.)} Let $b > 0$ and $\tau > 0$, and put
\[
  G(r) := b\int_0^r \bigl(1 - \Omega_d(\tau u)\bigr)\,\frac{du}{u} \qquad (r \ge 0) .
\]
Then $\xi \mapsto G(|\xi|)$ is the exponent of the datum $(0, b\,\mathbf 1_{(0,\tau)})$ in
dimension $d$, so that $G$ is the function of Proposition~\ref{prop:profile-one-dimension-down}(2), and in
Proposition~\ref{prop:radial-dictionary} it is the function of $c = b$ and of $X$ uniformly distributed on
the sphere of radius $\tau$. Hence $\zeta \mapsto \Omega_d(\tau|\zeta|)$ is not a characteristic
function on $\RR^{d+1}$.
\end{enumerate}
\end{corollary}

% A sentence of paper prose stands between the statement and its proof: `linkage check` reads a
% proof that directly follows a statement as that statement's proof, heading included.
Both clauses are read off the proposition through clause (1) of the dictionary.

\begin{proof}[Proof of Corollary~\ref{cor:dictionary-filtration}]
\emph{Clause (1).} Let $F$ and $X$ be as in
Proposition~\ref{prop:radial-dictionary}(1), with $\mathbb{E}\log^+|X| < \infty$, and
$F \in \mathcal{A}_{d+1}$. By that clause
$(0, c\,\Pr(|X| > \cdot))$ is a datum in dimension $d$ with exponent $F(|\cdot|)$, so its profile
is continuous, by Proposition~\ref{prop:profile-one-dimension-down}(1), and
$\Pr(|X| = r) = \Pr(|X| \ge r) - \Pr(|X| > r)$ is $1/c$ times the jump of
that profile at $r$, which is $0$.

\emph{Clause (2).} Let $U$ be uniformly distributed on $S^{d-1}$ and $X := \tau U$. It is
rotation invariant, $\Pr(X = 0) = 0$, $\mathbb{E}\log^+|X| = \log^+\tau < \infty$, and its radial
function is $\phi(r) = \Omega_d(\tau r)$. By Proposition~\ref{prop:radial-dictionary}(1) with $c = b$, the
function $G$ belongs to $\mathcal{A}_d$, and $\xi \mapsto G(|\xi|)$ is the exponent of the datum
$(0,k_d)$ with $k_d(r) = b\,\Pr(\tau > r) = b\,\mathbf 1_{(0,\tau)}(r)$. So $G$ is the function
of Proposition~\ref{prop:profile-one-dimension-down}(2), and $G \notin \mathcal{A}_{d+1}$ by that clause. By
Proposition~\ref{prop:radial-dictionary}(3), $G \notin \mathcal{A}_{d+1}$ says that
$\zeta \mapsto \phi(|\zeta|) = \Omega_d(\tau|\zeta|)$ is not a characteristic function on
$\RR^{d+1}$.
\end{proof}

% Q-0448: the proof of record of prop:step-exponent-dimensions (Q-0434), written by
% records/affine/transcribe-q0448-nodes.py and pinned. Its identity eq:shadow-identity is in
% the proof of Lemma lem:shadow-polar-profile, in the supplement.
\begin{proof}[Proof of Proposition~\ref{prop:step-exponent-dimensions}]
\emph{Four facts.} (a) For $n \ge 1$ and $\zeta \in \RR^n$, \ref{eq:sphere-mean-cos} with
$t = \tau$ and $x = \zeta$ says that $\zeta \mapsto \Omega_n(\tau|\zeta|)$ is the characteristic
function of the uniform probability measure on the sphere of radius $\tau$ of $\RR^n$, the image
of $\sigma_{n-1}/|S^{n-1}|$ under $\theta \mapsto \tau\theta$. (b) $\Omega_n$ is continuous and
$0 \le 1 - \Omega_n(t) \le t^2/2$, because $1 - \cos u \le u^2/2$ and
$|\langle e,\eta\rangle| \le 1$ on the sphere. (c) Let $(0,k)$ be an isotropic origin datum in
dimension $n$. By \ref{eq:polar-decomposition} with $g(x) = 1 - \cos\langle\xi,x\rangle$, by
Tonelli's theorem for a nonnegative integrand, and by the definition of $\Omega_n$ with
$e = \xi/|\xi|$, its exponent at $\xi \ne 0$ is
\[
  \frac{1}{|S^{n-1}|}\int_{S^{n-1}}\sigma_{n-1}(d\eta)\int_0^\infty
     \bigl(1 - \cos(r\langle\xi,\eta\rangle)\bigr)\,k(r)\,\frac{dr}{r}
   = \int_0^\infty \bigl(1 - \Omega_n(r|\xi|)\bigr)\,k(r)\,\frac{dr}{r} ,
\]
and both sides vanish at $\xi = 0$. For $k = b\,\mathbf 1_{(0,\tau)}$ this is $G_n(|\xi|)$, which
is finite by (b). (d) Let $\nu$ be a symmetric \Levy{} measure on $\RR^n$ and $v \ge 0$. The
infinitely divisible law with triplet $(vI, \nu, 0)$ that \ledger{A3}(iii) supplies has the
transform $\exp[-\tfrac12 v|\zeta|^2 - \int(1 - \cos\langle\zeta,z\rangle)\,\nu(dz)]$. Indeed
the integrand $e^{i\langle\zeta,z\rangle} - 1 - i\langle\zeta,z\rangle\mathbf 1_D(z)$ of the
representation has the real part $\cos\langle\zeta,z\rangle - 1$, and its imaginary part is odd
in $z$ and integrable against $\nu$, being at most $|\zeta|^2|z|^2/2$ in modulus on $D$ (as
$|\sin u - u| \le u^2/2$ for every real $u$) and at most $1$ off $D$; so the integral of the
imaginary part against the symmetric measure $\nu$ is $0$.

\emph{The shadow of the step two dimensions down.} The step datum
$(0, b\,\mathbf 1_{(0,\tau)})$ of dimension $d + 2$ has the exponent
$\zeta \mapsto G_{d+2}(|\zeta|)$, by (c). By Lemma~\ref{lem:shadow-polar-profile} with $d + 1$ in place
of $d$, its shadow $(0,k_1)$ is an isotropic origin datum in dimension $d + 1$ with exponent
$\zeta \mapsto G_{d+2}(|\zeta|)$, $\zeta \in \RR^{d+1}$, and
\[
  k_1(s) = C_{d+1}\int_{\RR} b\,\mathbf 1_{\{s\sqrt{1+w^2} < \tau\}}\,(1+w^2)^{-(d+2)/2}\,dw
         = 2b\,C_{d+1}\,J_{d+2}\Bigl(\frac{\tau}{s}\Bigr) \quad (0 < s < \tau),
\]
by the substitution $v = \sqrt{1+w^2}$ on $w > 0$, for which $dw = v\,(v^2-1)^{-1/2}dv$; and
$k_1(s) = 0$ for $s \ge \tau$. By the same lemma, the shadow $(0,k_2)$ of $(0,k_1)$ is an
isotropic origin datum in dimension $d$ with exponent $\xi \mapsto G_{d+2}(|\xi|)$,
$\xi \in \RR^d$, and
\[
  k_2(r) = C_d\,C_{d+1}\,b\int_{\RR}\int_{\RR}
     \mathbf 1_{\{r\sqrt{(1+u^2)(1+w^2)} < \tau\}}\,(1+w^2)^{-(d+2)/2}\,(1+u^2)^{-(d+1)/2}\,dw\,du .
\]
Substitute $w = y/\sqrt{1+u^2}$ in the inner integral. Then $(1+u^2)(1+w^2) = 1 + u^2 + y^2$ and
$(1+w^2)^{-(d+2)/2}(1+u^2)^{-(d+1)/2}\,dw = (1+u^2+y^2)^{-(d+2)/2}\,dy$, so with $q = (u,y)$
\[
  k_2(r) = C_d\,C_{d+1}\,b\int_{\RR^2}\mathbf 1_{\{r\sqrt{1+|q|^2} < \tau\}}\,
     (1+|q|^2)^{-(d+2)/2}\,dq .
\]
The same computation with the indicator replaced by $1$, and the normalisation of
Lemma~\ref{lem:shadow-polar-profile}(1) in the dimensions $d$ and $d + 1$, give
$C_d\,C_{d+1}\int_{\RR^2}(1+|q|^2)^{-(d+2)/2}\,dq = 1$. In polar coordinates of the plane, with
$v = \sqrt{1+|q|^2}$, the integral of $\mathbf 1_{\{v < x\}}(1+|q|^2)^{-(d+2)/2}$ is
$2\pi\int_1^x v^{-d-1}\,dv = (2\pi/d)(1 - x^{-d})$ for $x \ge 1$, and $0$ for $x < 1$. At
$x = \infty$ this gives
\begin{equation}\label{eq:shadow-constants}
  2\pi\,C_d\,C_{d+1} = d ,
\end{equation}
and at $x = \tau/r$ it gives
\[
  k_2(r) = b\Bigl(1 - \Bigl(\frac{r}{\tau}\Bigr)^{d}\Bigr) \quad (0 < r < \tau), \qquad
  k_2(r) = 0 \quad (r \ge \tau) .
\]

\emph{Clause (2), the identity.} Fix $\rho \ge 0$ and put $h(r) := 1 - \Omega_d(r\rho)$ and
$h_2(r) := 1 - \Omega_{d+2}(r\rho)$. By (c) for the datum $(0,k_2)$, whose exponent is
$G_{d+2}(|\cdot|)$, and by the definition of $G_{d+2}$,
\begin{equation}\label{eq:step-two-forms}
  \int_0^\tau h_2(r)\,\frac{dr}{r} = \frac{G_{d+2}(\rho)}{b}
   = \int_0^\tau h(r)\,\frac{dr}{r} - \tau^{-d}\int_0^\tau h(r)\,r^{d-1}\,dr .
\end{equation}
This holds for every $\tau > 0$, the parameter $\tau$ having been arbitrary. By (b) the
functions $h$ and $h_2$ are continuous with $0 \le h(r), h_2(r) \le r^2\rho^2/2$, so the three
integrals are continuously differentiable functions of $\tau > 0$. Differentiating
\ref{eq:step-two-forms} in $\tau$,
\[
  \frac{h_2(\tau)}{\tau} = \frac{h(\tau)}{\tau} + d\,\tau^{-d-1}\int_0^\tau h(r)\,r^{d-1}\,dr
     - \frac{h(\tau)}{\tau} ,
\]
that is
\[
  \tau^{-d}\int_0^\tau h(r)\,r^{d-1}\,dr = \frac{1}{d}\bigl(1 - \Omega_{d+2}(\tau\rho)\bigr) .
\]
By the definition of $G_d$ and \ref{eq:step-two-forms},
$G_d(\rho) - G_{d+2}(\rho) = b\,\tau^{-d}\int_0^\tau h(r)\,r^{d-1}\,dr$, which is
\ref{eq:step-exponent-two-up}.

\emph{Clause (2), the law.} Let $\mu^{\mathrm{step}}$ be the law of the step datum of dimension
$d + 2$, the infinitely divisible law with triplet $(0,\nu^{\mathrm{step}}_{d+2},0)$ that
\ledger{A3}(iii) supplies; its transform is $e^{-G_{d+2}(|\zeta|)}$, by (c) and (d). Let
$\pi := e^{-b/d}\sum_{n \ge 0}(b/d)^n\sigma_\tau^{*n}/n!$. It is a probability measure, and by
(a) its transform is
$e^{-b/d}\sum_n\bigl((b/d)\,\Omega_{d+2}(\tau|\zeta|)\bigr)^n/n!
= \exp\bigl[-\tfrac{b}{d}(1 - \Omega_{d+2}(\tau|\zeta|))\bigr]$. It is infinitely divisible: for
$m \ge 1$ the law defined in the same way with $b/(dm)$ in place of $b/d$ has the transform
$\widehat\pi^{1/m}$, so its $m$-fold convolution is $\pi$. The convolution of two infinitely
divisible laws is infinitely divisible, the convolution of their roots being a root. So
$\mu^{\mathrm{step}} * \pi$ is infinitely divisible, and by \ref{eq:step-exponent-two-up} its
transform is $e^{-G_d(|\zeta|)}$. Put
$\nu^{(d+2)} := \nu^{\mathrm{step}}_{d+2} + \tfrac{b}{d}\,\sigma_\tau$. It is the sum of a
\Levy{} measure and a finite measure that gives the origin no mass, so it is a \Levy{} measure,
and it is symmetric. By (a) and (c),
$\int(1 - \cos\langle\zeta,z\rangle)\,\nu^{(d+2)}(dz) = G_{d+2}(|\zeta|)
+ \tfrac{b}{d}(1 - \Omega_{d+2}(\tau|\zeta|)) = G_d(|\zeta|)$. The infinitely divisible law with
triplet $(0,\nu^{(d+2)},0)$, which \ledger{A3}(iii) supplies, therefore has the transform
$e^{-G_d(|\zeta|)}$, by (d). So it is $\mu^{\mathrm{step}} * \pi$, and the triplet of a law is
unique (\ledger{A3}(ii)).

\emph{Clause (3).} Write $k' = k_1 + k_2'$ on $(0,\infty)$, with $k_1$ as above and
\[
  k_2'(s) := (2b\,C_{d+1}/d)\,(s/\tau)^d\,s\,(\tau^2 - s^2)^{-1/2} \quad (0 < s < \tau), \qquad
  k_2'(s) := 0 \quad (s \ge \tau) .
\]
Both are nonnegative and Borel. The identity
\ref{eq:shadow-identity} holds for every nonnegative Borel profile $h$, monotone or not. For
$h = k_1$ its function $\bar h$ is $k_2$. For $h = k_2'$ and $0 < r < \tau$ put
$A := \sqrt{\tau^2 - r^2}$. The integrand of $\bar h(r)$ vanishes unless $|u| < A/r$, and there
$k_2'(r\sqrt{1+u^2})\,(1+u^2)^{-(d+1)/2} = (2b\,C_{d+1}/d)\,r^{d+1}\tau^{-d}\,(A^2 - r^2u^2)^{-1/2}$.
Since $\int_{-A/r}^{A/r}(A^2 - r^2u^2)^{-1/2}\,du = \pi/r$, and by \ref{eq:shadow-constants},
\[
  \bar h(r) = C_d\cdot\frac{2b\,C_{d+1}}{d}\cdot\frac{r^{d+1}}{\tau^{d}}\cdot\frac{\pi}{r}
            = b\Bigl(\frac{r}{\tau}\Bigr)^{d} \quad (0 < r < \tau), \qquad
  \bar h(r) = 0 \quad (r \ge \tau) .
\]
Adding the two, \ref{eq:shadow-identity} for $h = k'$ has $\bar h = k_2 + b\,(r/\tau)^d
\mathbf 1_{(0,\tau)} = b\,\mathbf 1_{(0,\tau)}$: for every nonnegative Borel $g$ on $\RR^d$
vanishing at the origin,
\begin{equation}\label{eq:step-one-up-image}
  \int_{\RR^{d+1}} g(\Pi z)\,\nu'(dz) = \int_{\RR^d} g\,d\nu^{\mathrm{step}}_d ,
\end{equation}
with $\Pi$ the projection on the first $d$ coordinates and $\nu^{\mathrm{step}}_d$ the \Levy{}
measure of the step datum of dimension $d$. The measure $\nu'$ is rotation invariant, having a
radial density, and gives the origin no mass. It is a \Levy{} measure: in polar coordinates of
$\RR^{d+1}$, $\int(|z|^2 \wedge 1)\,\nu'(dz) \le \int_0^\tau k'(s)\,s\,ds$, and $k'$ is bounded
on $(0,\tau/2]$ and at most a constant times $(\tau - s)^{-1/2}$ on $[\tau/2,\tau)$. Put
$\psi(\zeta) := \int(1 - \cos\langle\zeta,z\rangle)\,\nu'(dz)$ for $\zeta \in \RR^{d+1}$, which
is finite. It depends on $|\zeta|$ only, by the invariance of $\nu'$. For $\xi \in \RR^d$,
\ref{eq:step-one-up-image} with $g(x) = 1 - \cos\langle\xi,x\rangle$ gives
$\psi((\xi,0)) = G_d(|\xi|)$, by (c). So $\psi(\zeta) = G_d(|\zeta|)$ for every
$\zeta \in \RR^{d+1}$. The infinitely divisible law with triplet $(0,\nu',0)$, which
\ledger{A3}(iii) supplies, has the transform $e^{-\psi}$, by (d), and its triplet is unique (\ledger{A3}(ii)). The properties of $k'$: both terms are
positive on $(0,\tau)$, since $J_{d+2}(x) > 0$ for $x > 1$. As $s \downarrow 0$ the second term
tends to $0$, and $2C_{d+1}J_{d+2}(\tau/s) \to 2C_{d+1}J_{d+2}(\infty)
= C_{d+1}\int_{\RR}(1+w^2)^{-(d+2)/2}\,dw = 1$, by Lemma~\ref{lem:shadow-polar-profile}(1) in dimension
$d + 1$. As $s \uparrow \tau$ the first term tends to $0$ and the second to $\infty$. The closed
forms. By Lemma~\ref{lem:shadow-polar-profile}(1), $C_2 = (\int_{\RR}(1+u^2)^{-3/2}du)^{-1} = \tfrac12$
and $C_3 = (\int_{\RR}(1+u^2)^{-2}du)^{-1} = 2/\pi$. Differentiation shows
$J_3(x) = \sqrt{x^2-1}/x$ and $J_4(x) = \tfrac12\bigl[\sqrt{x^2-1}/x^2 + \arccos(1/x)\bigr]$,
both sides vanishing at $x = 1$. For $d = 1$ this gives
$k'(s) = b\bigl[\sqrt{\tau^2-s^2}/\tau + s^2/(\tau\sqrt{\tau^2-s^2})\bigr]
= b\,\tau/\sqrt{\tau^2-s^2}$, and for $d = 2$,
$k'(s) = (2b/\pi)\bigl[s\sqrt{\tau^2-s^2}/\tau^2 + \arccos(s/\tau)
+ s^3/(\tau^2\sqrt{\tau^2-s^2})\bigr] = (2b/\pi)\bigl[\arccos(s/\tau) + s/\sqrt{\tau^2-s^2}\bigr]$.

\emph{Clause (1), ``if''.} Let $D \le d + 2$, let $\mu$ be the law of clause (2), and let $P$ be
the projection of $\RR^{d+2}$ on its first $D$ coordinates. The image $P_*\mu$ has the transform
$\zeta \mapsto \widehat\mu((\zeta,0)) = e^{-G_d(|\zeta|)}$, $\zeta \in \RR^D$. It is infinitely
divisible: the image of a convolution under a linear map is the convolution of the images, so
the image of an $m$-th convolution root of $\mu$ is one of $P_*\mu$.

\emph{Clause (1), ``only if''.} Let $D \ge d + 3$ and suppose that an infinitely divisible law on
$\RR^D$ has the transform $e^{-G_d(|\zeta|)}$. Its image under the projection on the first
$d + 3$ coordinates is infinitely divisible with the transform $e^{-G_d(|\zeta|)}$ on
$\RR^{d+3}$, as in the previous paragraph. So let $D = d + 3$, and let $\mu'$ be such a law. By
\ledger{A3}(i) it has a triplet $(A,\nu',\gamma)$. Its transform is invariant under $\Orth(D)$,
so, as in the first paragraph of the proof of Proposition~\ref{prop:isotropic-polar-reading} from its
third sentence on, which reads of
the law only that it is infinitely divisible with a radial transform, \ledger{A3}(ii) gives
$\langle\zeta,A\zeta\rangle = \alpha|\zeta|^2$ with $\alpha \ge 0$ and $Q_*\nu' = \nu'$ for every
$Q \in \Orth(D)$, and
\[
  G_d(|\zeta|) = \tfrac12\alpha|\zeta|^2
     + \int_{\RR^{d+3}}\bigl(1 - \cos\langle\zeta,z\rangle\bigr)\,\nu'(dz), \qquad \zeta \in \RR^{d+3} .
\]
Only the modulus of the representation is used for this display: $\mu'$ is not assumed
symmetric, and $-\log|\widehat{\mu'}(\zeta)| = G_d(|\zeta|)$ because $e^{-G_d}$ is real and
positive.
Let $\Pi$ be the projection of $\RR^{d+3}$ on its first $d + 2$ coordinates, and let $\bar\nu$ be
the image of $\nu'$ under $\Pi$, restricted to $\RR^{d+2} \setminus \{0\}$. It is symmetric, and
it is a \Levy{} measure, because $|\Pi z| \le |z|$ gives
$\int(|y|^2 \wedge 1)\,\bar\nu(dy) \le \int(|z|^2 \wedge 1)\,\nu'(dz) < \infty$. For
$\eta \in \RR^{d+2}$, $\langle(\eta,0),z\rangle = \langle\eta,\Pi z\rangle$, and the integrand
vanishes where $\Pi z = 0$, so
$G_d(|\eta|) = \tfrac12\alpha|\eta|^2 + \int(1 - \cos\langle\eta,y\rangle)\,\bar\nu(dy)$. By
\ledger{A3}(iii) there is an infinitely divisible law on $\RR^{d+2}$ with triplet
$(\alpha I,\bar\nu,0)$, and by (d) its transform is $e^{-G_d(|\eta|)}$. The law of clause (2) has the
same transform and the triplet $(0,\nu^{(d+2)},0)$. The two laws are equal, and by
\ledger{A3}(ii) $\alpha = 0$ and $\bar\nu = \nu^{(d+2)}$. Let
$E := \{y \in \RR^{d+2} : |y| = \tau\}$. Then
$\bar\nu(E) = \nu^{(d+2)}(E) \ge \tfrac{b}{d}\,\sigma_\tau(E) = \tfrac{b}{d} > 0$.

On the other hand $\bar\nu(E) = \nu'(\Pi^{-1}E)$, and
$\Pi^{-1}E = \{z \in \RR^{d+3} : |\Pi z| = \tau\}$ lies in $\{|z| \ge \tau\}$. Let $m$ be the
restriction of $\nu'$ to $\{|z| \ge \tau\}$. It is a finite measure, because
$\int(|z|^2 \wedge 1)\,d\nu' < \infty$, it is invariant under $\Orth(d+3)$, and
$m(\Pi^{-1}E) = \bar\nu(E) > 0$, so $m \ne 0$. Let $Z$ be a random vector in $\RR^{d+3}$ with the
law $m/m(\RR^{d+3})$. It is rotation invariant, so by the preamble of
Proposition~\ref{prop:radial-dictionary} it has the law of $R\,U$, with $R$ of the law of $|Z|$, so that
$R \ge \tau$, and $U$ uniformly distributed on $S^{d+2}$ and independent of $R$. Hence
\[
  \frac{m(\Pi^{-1}E)}{m(\RR^{d+3})} = \Pr\bigl(R\,|\Pi U| = \tau\bigr)
   = \int \Pr\bigl(|\Pi U| = \tau/r\bigr)\,\Pr(R \in dr) .
\]
For every $w \in [0,1]$, $\Pr(|\Pi U| = w) = 0$. Indeed $|\Pi u|^2 = 1 - u_{d+3}^2$ on the
sphere, so the event is $\{U \in K\}$ with $K := \{z \in \RR^{d+3} : z_{d+3}^2 = (1-w^2)|z|^2\}$,
a cone. The cone is Lebesgue-null: if $w > 0$ it is the set
$\{w^2z_{d+3}^2 = (1-w^2)|\Pi z|^2\}$, which meets every line parallel to the last axis in at
most two points, and if $w = 0$ it is the last axis. In polar coordinates the Lebesgue measure
of $K \cap \{1 < |z| < 2\}$ is $\sigma_{d+2}(K \cap S^{d+2})\int_1^2 r^{d+2}\,dr$, so
$\sigma_{d+2}(K \cap S^{d+2}) = 0$. Therefore $m(\Pi^{-1}E) = 0$, a contradiction.

\emph{The last sentence.} For $d = 1$, $G(r) = b\,\Cin(\tau r)$ and the ray is in no
$\mathcal{A}_D$ with $D \ge 2$, by Proposition~\ref{prop:profile-one-dimension-down}(3). The rest is clause
(1) with $d + 2 = 3$.
\end{proof}

The last proposition says what the members of the separating family of
Definition~\ref{def:filtration-thresholds} are as laws.

% shared with blueprint prop:separating-family-mixtures+5fd15a56009e
\begin{proposition}[the separating family as laws]\label{prop:separating-family-mixtures}
Let $w > 0$, and let $F^{(w)}_{\mathrm I}$, $F^{(w)}$ and the critical values $w_D$ be as in
Definition~\ref{def:filtration-thresholds}.
\begin{enumerate}
\item \emph{(Closed form.)} For $\sigma \ge 0$ and $r \ge 0$,
\[
  F^{(w)}_{\mathrm I}(\sigma) = 2\sqrt{\pi\sigma} + w\int_1^2\bigl(1 - e^{-\sigma u}\bigr)\frac{du}{u},
  \qquad
  F^{(w)}(r) = \sqrt{2\pi}\,r + w\int_1^2\bigl(1 - e^{-u r^2/2}\bigr)\frac{du}{u} .
\]
\item \emph{(A Gaussian variance mixture in every dimension.)} Let $Z_0$ be a standard Gaussian
real variable, $N$ a Poisson variable of mean $w\log 2$, and $U_1, U_2, \dots$ variables with the
density $(u\log 2)^{-1}$ on $[1,2]$, all independent, and put
\[
  W := \frac{2\pi}{Z_0^2} + \sum_{j=1}^{N}U_j .
\]
Then $W > 0$ almost surely and $\mathbb{E}\,e^{-\sigma W} = e^{-F^{(w)}_{\mathrm I}(\sigma)}$ for
every $\sigma \ge 0$. For every integer $D \ge 1$, if $Z$ is a standard Gaussian vector of
$\RR^D$ independent of $W$, then $\sqrt W\,Z$ has the characteristic function
$\xi \mapsto e^{-F^{(w)}(|\xi|)}$ on $\RR^D$. Its law $\mu^{(w)}_D$ is rotation invariant and
infinitely divisible. It is the convolution of the law with characteristic function
$e^{-\sqrt{2\pi}|\xi|}$, a rotation-invariant Cauchy law, and of a compound Poisson law, with
characteristic function $\exp\bigl[-w\int_1^2(1 - e^{-u|\xi|^2/2})\,du/u\bigr]$.
\item \emph{(Self-decomposable up to the critical dimension.)} For every integer $D \ge 1$,
$\mu^{(w)}_D$ is self-decomposable if and only if $F^{(w)} \in \mathcal{A}_D$, that is, by
Theorem~\ref{thm:separating-family-critical-values}(1), if and only if $w \le w_D$. Consequently, if
$w \in (w_{d+1}, w_d\,]$ for an integer
$d \ge 1$, then $\mu^{(w)}_D$ is self-decomposable exactly for $D \le d$, and if $w > w_1$ it is
self-decomposable in no dimension.
\end{enumerate}
\end{proposition}

% A sentence of paper prose stands between the statement and its proof.
In clause (2) the two summands of $W$ are a one-sided stable variable of index $\tfrac12$ and a
compound Poisson sum.

\begin{proof}[Proof of Proposition~\ref{prop:separating-family-mixtures}]
\emph{Clause (1).} By Definition~\ref{def:filtration-thresholds} the delay datum $D^{(w)}$ has drift $0$
and profile $k^{(w)}_{\mathrm I}(u) = u^{-1/2} + w\,\mathbf 1_{[1,2]}(u)$, so
$F^{(w)}_{\mathrm I}(\sigma) = \int_0^\infty(1 - e^{-\sigma u})\,u^{-3/2}\,du +
w\int_1^2(1 - e^{-\sigma u})\,du/u$. Integrating by parts, for $\sigma > 0$,
\[
\begin{aligned}
  \int_0^\infty\bigl(1 - e^{-\sigma u}\bigr)\,u^{-3/2}\,du
   &= \Bigl[-2u^{-1/2}\bigl(1 - e^{-\sigma u}\bigr)\Bigr]_0^\infty
     + 2\sigma\int_0^\infty u^{-1/2}e^{-\sigma u}\,du \\
   &= 2\sigma\cdot\Gamma(\tfrac12)\,\sigma^{-1/2} = 2\sqrt{\pi\sigma} ,
\end{aligned}
\]
the bracket vanishing at both ends. At $\sigma = 0$ both sides are $0$. The second formula is the
first at $\sigma = r^2/2$.

\emph{Clause (2), a Gaussian integral.} For $\beta \ge 0$,
\begin{equation}\label{eq:gaussian-inverse-square}
  \int_0^\infty e^{-x^2 - \beta^2/x^2}\,dx = \frac{\sqrt\pi}{2}\,e^{-2\beta} .
\end{equation}
For $\beta = 0$ this is the Gaussian integral. Let $\beta > 0$. Then
$x^2 + \beta^2/x^2 = (x - \beta/x)^2 + 2\beta$. Put
$J := \int_0^\infty e^{-(x - \beta/x)^2}dx$. The substitution $x = \beta/y$ gives
$J = \int_0^\infty e^{-(y - \beta/y)^2}\,\beta y^{-2}\,dy$. Adding the two expressions,
$2J = \int_0^\infty e^{-(x - \beta/x)^2}(1 + \beta/x^2)\,dx = \int_{\RR}e^{-t^2}dt = \sqrt\pi$,
by the substitution $t = x - \beta/x$, an increasing bijection of $(0,\infty)$ onto $\RR$ with
$dt = (1 + \beta/x^2)\,dx$.

\emph{Clause (2), the Laplace transform of $W$.} Let $\varkappa > 0$ and $\sigma \ge 0$. With
$z = \sqrt2\,x$ and \ref{eq:gaussian-inverse-square} at $\beta^2 = \sigma\varkappa/2$,
\[
  \mathbb{E}\,e^{-\sigma\varkappa/Z_0^2}
   = \frac{2}{\sqrt{2\pi}}\int_0^\infty e^{-z^2/2 - \sigma\varkappa/z^2}\,dz
   = \frac{2}{\sqrt\pi}\int_0^\infty e^{-x^2 - \sigma\varkappa/(2x^2)}\,dx
   = e^{-\sqrt{2\sigma\varkappa}} .
\]
At $\varkappa = 2\pi$ this is $e^{-2\sqrt{\pi\sigma}}$. For the sum, conditioning on $N$,
\[
\begin{aligned}
  \mathbb{E}\,e^{-\sigma\sum_{j \le N}U_j}
   &= \sum_{n \ge 0}e^{-w\log 2}\frac{(w\log 2)^n}{n!}\bigl(\mathbb{E}\,e^{-\sigma U_1}\bigr)^n
   = \exp\bigl[-w\log 2\,\bigl(1 - \mathbb{E}\,e^{-\sigma U_1}\bigr)\bigr] \\
   &= \exp\Bigl[-w\int_1^2\bigl(1 - e^{-\sigma u}\bigr)\frac{du}{u}\Bigr] ,
\end{aligned}
\]
because $\log 2\,(1 - \mathbb{E}\,e^{-\sigma U_1}) = \int_1^2(1 - e^{-\sigma u})\,du/u$. By
independence and clause (1), $\mathbb{E}\,e^{-\sigma W} = e^{-F^{(w)}_{\mathrm I}(\sigma)}$. And
$W \ge 2\pi/Z_0^2 > 0$ almost surely.

\emph{Clause (2), the mixture.} Let $D \ge 1$ and let $Z$ be a standard Gaussian vector of
$\RR^D$ independent of $W$. Conditioning on $W$, and using that $\sqrt u\,Z$ has the
characteristic function $e^{-u|\xi|^2/2}$,
\[
  \mathbb{E}\,e^{i\langle\xi,\sqrt W Z\rangle} = \mathbb{E}\,e^{-W|\xi|^2/2}
   = e^{-F^{(w)}_{\mathrm I}(|\xi|^2/2)} = e^{-F^{(w)}(|\xi|)} .
\]
The characteristic function is radial, so the law $\mu^{(w)}_D$ of $\sqrt W\,Z$ is invariant
under $\Orth(D)$, by the uniqueness theorem for characteristic functions. By clause (1) it
factors as $e^{-\sqrt{2\pi}|\xi|}$ times
$\exp[-w\int_1^2(1 - e^{-u|\xi|^2/2})\,du/u]$. The same computation with $W$ replaced by each
of its two summands shows that the first factor is the characteristic function of
$\sqrt{2\pi}\,Z/|Z_0|$ and the second that of $\sqrt{\sum_{j\le N}U_j}\,Z$. The second is also
the characteristic function of the compound Poisson sum $\sum_{j \le N}\sqrt{U_j}\,Z_j$, with
$Z_1, Z_2, \dots$ independent standard Gaussian vectors independent of the rest, by the
computation made above for the Laplace transform. So $\mu^{(w)}_D$ is the convolution of the
two laws.

\emph{Clause (2), infinite divisibility.} Let $n \ge 1$, let $N_n$ be a Poisson variable of mean
$(w\log 2)/n$, independent of the rest, and put
$W_n := 2\pi/(n^2Z_0^2) + \sum_{j \le N_n}U_j$. The two computations above, at
$\varkappa = 2\pi/n^2$ and with $w/n$ in place of $w$, give
$\mathbb{E}\,e^{-\sigma W_n} = e^{-F^{(w)}_{\mathrm I}(\sigma)/n}$. So the law $\mu_n$ of
$\sqrt{W_n}\,Z$ has the characteristic function $e^{-F^{(w)}(|\xi|)/n}$, and its $n$-fold
convolution has the characteristic function of $\mu^{(w)}_D$. Hence
$\mu^{(w)}_D = \mu_n^{\ast n}$ for every $n$: it is infinitely divisible.

\emph{Clause (3), the reading in laws.} $F^{(w)} \in \mathcal{A}_D$ means that $F^{(w)}$ is
admissible and that $\xi \mapsto F^{(w)}(|\xi|)$ is the exponent of a self-decomposable law on
$\RR^D$. A law with the characteristic function $e^{-F^{(w)}(|\xi|)}$ is $\mu^{(w)}_D$, by
clause (2) and the uniqueness theorem. And if $\mu^{(w)}_D$ is self-decomposable, $F^{(w)}$ is
admissible, by Proposition~\ref{prop:isotropic-polar-reading}(2). So $F^{(w)} \in \mathcal{A}_D$
exactly when $\mu^{(w)}_D$ is self-decomposable.

\emph{Clause (3), the criterion and the windows.} By
Theorem~\ref{thm:separating-family-critical-values}(1), $F^{(w)} \in \mathcal{A}_D$ exactly when
$w \le w_D$. By its clause (4) this holds exactly for $D \le d$ when $w \in (w_{d+1}, w_d\,]$,
and for no $D$ when $w > w_1$. With the reading in laws these are the assertions of the clause.
\end{proof}

